%%%%%%%%%%%%%%%%%%%%%%%%%%%%%%%%%%%%%%%%%%%%%%%%%%%%%%%%%%%%%%%%%%%%%%
% Lớp 12 | Một số yếu tố xác suất | Công thức xác suất toàn phần. Công thức Bayes
% Dạng 2D6?2-3: Tính xác suất bằng công thức xác suất Bayes
% Số câu: 84  (sinh tự động bởi tach_cau.py; từ nay thầy sửa trực tiếp file này)
%%%%%%%%%%%%%%%%%%%%%%%%%%%%%%%%%%%%%%%%%%%%%%%%%%%%%%%%%%%%%%%%%%%%%%

%%%%%%%%%%%%%%%% TRẮC NGHIỆM 4 PHƯƠNG ÁN %%%%%%%%%%%%%%%%

% nguồn: Lop_12_Dai_so_Chuong_6
\begin{ex}%[2D6N2-3]%[ma:5f684f31]
	Cho hai biến cố $A$, $B$ thỏa mãn  $\mathrm{P}\left(A\right)=0{,}4$; $\mathrm{P} \left(B\right)=0{,}3$; $\mathrm{P}\left(A|B\right)=0{,}25$. Khi đó, $\mathrm{P}\left(B|A\right)$ bằng
	\choice 
		 {\True $0{,}1875$}
		{$0{,}48$} 
		{$0{,}333$} 
		{$0{,}95$} 
	 \loigiai{
		Theo công thức Bayes, ta có 
	\[\mathrm{P}\left(B|A\right)=\dfrac{\mathrm{P} \left(B\right) \cdot \mathrm{P}\left(A|B\right)}{\mathrm{P}\left(A\right)}=\dfrac{0{,}3\cdot 0{,}25}{0{,}4}=0{,}1875.\]} 
\end{ex}

% nguồn: Lop_12_Dai_so_Chuong_6
\begin{ex}%[2D6N2-3]%[ma:f56e5205]
	Cho hai biến cố $A$ và $B$ sao cho $\mathrm{P}\left(A\right)=0{,}6$; $\mathrm{P}\left(B\right)=0{,}4$; $\mathrm{P}\left(A|B\right)=0{,}3$. Khi đó $\mathrm{P}\left(B|A\right)$ bằng bao  nhiêu?
	\choice 
		 {\True $0{,}2$}
		 {$0{,}3$}
		 {$0{,}4$}
		 {$0{,}6$}
	 \loigiai{
	Áp dụng công thức  Bayes, ta có \[\mathrm{P} \left(B|A\right)=\dfrac{\mathrm{P} \left(B\right)\mathrm{P} \left(A|B\right)}{\mathrm{P}\left(A\right)}=\dfrac{0{,}4 \cdot 0{,}3}{0{,}6}=0{,}2.\]}
\end{ex}

% nguồn: Lop_12_Dai_so_Chuong_6
\begin{ex}%[2D6N2-3]%[ma:41665ae4]
	Cho hai biến cố  $A$ và $B$, với $\mathrm{P} \left(A\right)=0{,}2$; $\mathrm{P} \left(B\right)=0{,}26$; $\mathrm{P}\left(B|A\right)=0{,}7$. Tính $\mathrm{P}\left(A|B\right)$.
	 \choice
		 {\True $\dfrac{7}{13}$}
		{$\dfrac{6}{13}$} 
		 {$\dfrac{4}{13}$}
		{$\dfrac{9}{13}$} 
	\loigiai{ 
	Theo công thức Bayes, ta có \[\mathrm{P} \left(A|B\right)=\dfrac{\mathrm{P} \left(A\right) \cdot  \mathrm{P} \left(B|A\right)}{\mathrm{P} \left(B\right)}=\dfrac{0{,}2 \cdot 0{,}7}{0{,}26}=\dfrac{7}{13}.\]}
\end{ex}

% nguồn: Lop_12_Dai_so_Chuong_6
\begin{ex}%[2D6N2-3]%[ma:84f140d9]
	Cho $\mathrm{P}(A)=0{,}3$; $\mathrm{P}(B)=0{,}4$  và $\mathrm{P}(A|B)=0{,}6$. Tính $\mathrm{P}(B|A)$.
	 \choice
		{$\mathrm{P}(B|A)=0{,}45$} 
		{\True  $\mathrm{P}(B|A)=0{,}8$}
		{$\mathrm{P}(B|A)=0{,}2$} 
		{$\mathrm{P}(B|A)=0{,}3$} 
	 \loigiai{
		 Áp dụng công thức Bayes ta có
		 \[\mathrm{P}(B|A)=\dfrac{\mathrm{P}(A|B) \cdot \mathrm{P}(B)}{\mathrm{P}(A)}=\dfrac{0{,}6 \cdot 0{,}4}{0{,}3}=0{,}8.\]
	 }
\end{ex}

% nguồn: Lop_12_Dai_so_Chuong_6
\begin{ex}%[2D6N2-3]%[ma:469e4a77]
	Cho hai  biến cố $A$ và $B$, với $\mathrm{P} \left(B\right)=0{,}8$; $\mathrm{P} \left(A|B\right)=0{,}7$; $\mathrm{P} \left(A|\overline{B}\right)=0{,}45$. Tính $\mathrm{P}\left(B|A\right)$.
	 \choice
		{$0{,}25$} 
		{\True  $\dfrac{56}{65}$}
		 {$0{,}65$}
		{$0{,}5$} 
	 \loigiai{
		Ta  có $\mathrm{P} \left(\overline{B}\right)=1-0{,}8=0{,}2$. \\
		Công thức Bayes 
		\[\mathrm{P} \left(B|A\right)=\dfrac{\mathrm{P}  \left(B\right)\mathrm{P} \left(A|B\right)}{\mathrm{P} \left(B\right)\mathrm{P} \left(A|B \right) +\mathrm{P} \left(\overline{B}\right)\mathrm{P} \left(A|\overline{B}\right)}
	=\dfrac{0{,}8  \cdot 0{,}7}{0{,}8 \cdot 0{,}7+0{,}2 \cdot 0{,}45}=\dfrac{56}{65}.\]}
\end{ex}

% nguồn: Lop_12_Dai_so_Chuong_6
\begin{ex}%[2D6N2-3]%[ma:954ada76]
	Cho hai biến cố $A$  và $B$, với $\mathrm{P}\left(A\right)=0{,}2$; $\mathrm{P} \left(B|A\right)=0{,}7$; $\mathrm{P} \left(B|\overline{A}\right)=0{,}15$. Tính $\mathrm{P} \left(A|B\right)$.
	\choice 
		{\True $\dfrac{7}{13}$} 
		 {$\dfrac{6}{13}$}
		{$\dfrac{4}{13}$} 
		{$\dfrac{9}{13}$} 
	 \loigiai{
		 Ta có $\mathrm{P} \left(A\right)=0{,}2 \Rightarrow \mathrm{P} \left(\overline{A}\right)=0{,}8$; $\mathrm{P}\left(B|A\right)=0{,}7$; $\mathrm{P}\left(B|\overline{A}\right)=0{,}15$.\\
		 Ta có
		\[\mathrm{P} \left(B\right)=\mathrm{P} \left(A\right) \cdot \mathrm{P} \left(B|A\right)+\mathrm{P} \left(\overline{A}\right) \cdot  \mathrm{P} \left(B|\overline{A}\right) = \mathrm{P}\left(B\right)=0{,}2 \cdot 0{,}7+0{,}8 \cdot 0{,}15=0{,}26.\]
		Theo công  thức Bayes, ta có
	\[\mathrm{P} \left(A|B\right)=\dfrac{\mathrm{P} \left(A\right) \cdot \mathrm{P}\left(B|A\right)}{\mathrm{P}\left(B\right)}  =\dfrac{0{,}2 \cdot 0{,}7}{0{,}26}=\dfrac{7}{13}.\]}
\end{ex}

% nguồn: Lop_12_Dai_so_Chuong_6
\begin{ex}%[2D6H2-3]%[ma:f974b49f]
	Một bệnh viện sử dụng một xét nghiệm để phát hiện một loại bệnh với độ chính xác là $95\%$ (nghĩa là $95\%$ bệnh nhân mắc bệnh sẽ  có kết quả dương tính). Xét nghiệm này cũng có tỷ lệ dương tính già là $2\%$ (nghĩa là $2\%$ bệnh nhân không mắc bệnh cũng có kết quả dương tính). Biết rằng $1\%$ dân số thực sự mắc bệnh này. Nếu một người nhận kết quả xét nghiệm dương tính, xác suất thực sự người đó mắc bệnh là bao nhiêu?
	 \choice
		 {\True Khoảng $32\%$}
		{Khoảng $47\%$} 
		{Khoảng  $83\%$}
		{Khoảng  $95\%$}
	\loigiai{ 
		Gọi $B$ là biến cố \lq\lq  người mắc bệnh\rq\rq, có $\mathrm{P}(B)=1\%=0{,}01$.\\
		Suy ra $\mathrm{P}\left(\overline{B}\right)=1-\mathrm{P}(B)=1-0{,}01=0{,}99$.\\ 
		Gọi $D$ là biến cố \lq\lq người đó được xét nghiệm có kết  quả dương tính\rq\rq.\\
		Khi đó, xác suất xét nghiệm  dương tính nếu mắc bệnh là $\mathrm{P}(D\mid B)=95\%=0{,}95$.\\
		Xác suất xét nghiệm dương tính  nếu không mắc bệnh là $\mathrm{P}\left(D\mid \overline{B}\right)=2\%=0{,}02$.\\
		Dùng công thức Bayes để tính xác suất mắc bệnh  khi có kết quả xét nghiệm dương tính
		\allowdisplaybreaks 
		\begin{align*}
			\mathrm{P}(B\mid D)&=\dfrac{\mathrm{P}(B)\cdot\mathrm{P}(D\mid B)}{\mathrm{P}(B)\cdot\mathrm{P}(D\mid B)+\mathrm{P}\left(\overline{B}\right)\cdot\mathrm{P}\left(D\mid  \overline{B}\right)}\\
			&=\dfrac{0{,}95  \cdot 0{,}01}{0{,}95 \cdot 0{,}01+0{,}02 \cdot 0{,}99}\\
			&\approx\ 0{,}3242. 
		\end{align*}
		Vậy xác suất người đó thực sự mắc  bệnh là khoảng $32\%$.
	} 
\end{ex}

% nguồn: Lop_12_Dai_so_Chuong_6
\begin{ex}%[2D6H2-3]%[ma:a5dd4171]
	Được biết có $5\%$ đàn ông bị mù màu, và $0{,}25\%$ phụ nữ bị mù màu (Nguồn: F. M. Dekking et al., A modern introduction to probability and statistics-Understanding why and how, Springer, 2005).  Giả sử số đàn ông bằng số phụ nữ. Chon một người bị mù màu. Xác suất để người đó là đàn ông là bao nhiêu?
	 \choice
		 {$\dfrac{19}{21}$}
		 {\True $\dfrac{20}{21}$}
		{$\dfrac{24}{25}$} 
		{$\dfrac{18}{25}$} 
	\loigiai{ 
		Gọi $A$ là biến cố \lq\lq người được chọn là đàn ông\rq\rq, $B$ là biến cố \lq\lq người được  chọn mù màu\rq\rq.\\
		Theo đề bài ra ta có  $\mathrm{P}(B\mid A)=0{,}05$; $\mathrm{P}\left(B\mid \overline{A}\right)=0{,}0025$.\\
		Vì số đàn ông bằng số phụ nữ nên ta có $\mathrm{P}(A)=\mathrm{P}\left(\overline{A}\right)=0{,}5$.\\ 
		Áp dụng công thức Bayes ta có xác suất để chọn được một người đàn ông mù  màu là
		$$ 
		\mathrm{P}(A\mid B)=\dfrac{\mathrm{P}(A) \cdot \mathrm{P}(B\mid A)}{\mathrm{P}(A) \cdot \mathrm{P}(B\mid A)+\mathrm{P}\left(\overline{A}\right) \cdot \mathrm{P}\left(B\mid \overline{A}\right)}=\dfrac{0{,}5 \cdot 0{,}05}{0{,}5 \cdot  0{,}05+0{,}5 \cdot 0{,}0025}=\dfrac{20}{21}.
		 $$
	 }
\end{ex}

% nguồn: Lop_12_Dai_so_Chuong_6
\begin{ex}%[2D6H2-3]%[ma:41f73f04]
	Có hai đội thi đấu môn bơi lội. Đội $I$ có $4$ vận động viên, đội $II$ có $6$ vận động viên. Xác suất đạt huy chương bạc của mỗi vận động viên đội $I$ và đội $II$ tương ứng là $0{,}7$ và $0{,}6$. Chọn ngẫu nhiên một vận động viên. Giả sử vận động viên được chọn  đạt huy chương bạc. Tính xác suất để vận động viên này thuộc đội $I$.
	 \choice
		{$\dfrac{8}{11}$} 
		{$\dfrac{11}{16}$} 
		 {$\dfrac{3}{16}$}
		{\True $\dfrac{7}{16}$} 
	 \loigiai{
		Gọi $A$ là biến cố: \lq\lq Vận động viên được chọn thuộc đội $I$  \rq\rq.\\
		Suy ra $\overline{A}$ là biến cố: \lq\lq Vận động viên được chọn thuộc  đội $II$ \rq\rq.\\
		$B$ là biến cố: \lq\lq Vận động  viên được chọn đạt huy chương bạc\rq\rq.\\
		Khi đó $\mathrm{P}(B)=\mathrm{P}(A) \cdot  \mathrm{P}(B\mid A)+\mathrm{P}\left(\overline{A}\right) \cdot \mathrm{P}\left(B\mid \overline{A}\right)$.\\
		Trong đó: $\mathrm{P}(A)=\dfrac{4}{10}=\dfrac{2}{5}$, $\mathrm{P}(B\mid A)=0{,}7$;  $\mathrm{P}\left(\overline{A}\right)=\dfrac{6}{10}=\dfrac{3}{5}$, $\mathrm{P}\left(B\mid \overline{A}\right)=0{,}6$.\\
		Suy  ra $\mathrm{P}(B)=\dfrac{2}{5} \cdot 0{,}7+\dfrac{3}{5} \cdot 0{,}6=\dfrac{16}{25}$.\\
		Theo công  thức Bayes ta có $\mathrm{P}(A\mid B)=\dfrac{\mathrm{P}(A) \cdot \mathrm{P}(B\mid A)}{\mathrm{P}(B)}=\dfrac{\dfrac{2}{5} \cdot 0{,}7}{\dfrac{16}{25}}=\dfrac{7}{16}$.
	 }
\end{ex}

% nguồn: Lop_12_Dai_so_Chuong_6
\begin{ex}%[2D6H2-3]%[ma:70a2fb13]
	Trường THPT Hòa Bình có $20\%$ học sinh tham gia câu lạc bộ âm nhạc, trong số học sinh đó có $85\%$ học sinh biết chơi đàn guitar. Ngoài ra, có $10\%$ số học sinh không tham gia câu lạc bộ âm nhạc  cũng biết chơi đàn guitar. Chọn ngẫu nhiên $1$ học sinh của trường. Giả sử học sinh đó biết chơi đàn guitar. Xác suất chọn được học sinh thuộc câu lạc bộ âm nhạc là
	\choice 
		{\True $\dfrac{17}{25}$} 
		{$\dfrac{7}{25}$} 
		{$\dfrac{17}{29}$} 
		 {$\dfrac{17}{75}$}
	\loigiai{ 
		Xét các biến cố: $A$:  \lq\lq Chọn được học sinh thuộc câu lạc bộ âm nhạc\rq\rq\;\\
		$B$:  \lq\lq Chọn được học sinh biết chơi đàn guitar\rq\rq.\\
		Khi đó, $\mathrm{P}(A)=0{,}2$; $\mathrm{P}\left(\overline{A}\right)=0{,}8$; $\mathrm{P}(B\mid A)=0{,}85$; $\mathrm{P}\left(B\mid  \overline{A}\right)=0{,}1$.\\
		Theo công thức  xác suất toàn phần, ta có
		 $$
		\mathrm{P}(B)=\mathrm{P}(A) \cdot \mathrm{P}(B \mid A)+\mathrm{P}\left(\overline{A}\right) \cdot \mathrm{P}\left(B \mid \overline{A}\right)=0{,}2  \cdot 0{,}85+0{,}8 \cdot 0{,}1=0{,}25.
		$$ 
		Theo công thức Bayes, ta có $\mathrm{P}(A\mid B)=\dfrac{\mathrm{P}(A)  \cdot \mathrm{P}(B\mid A)}{\mathrm{P}(B)}=\dfrac{0{,}2\cdot 0{,}85}{0{,}25}=0{,}68=\dfrac{17}{25}$.
	} 
\end{ex}

% nguồn: Lop_12_Dai_so_Chuong_6
\begin{ex}%[2D6H2-3]%[ma:65b56dea]
	Năm $2001$, Cộng đồng châu Âu có làm một đợt kiểm tra rất rộng rãi các con bò để phát hiện những con bò bị bệnh bò điên. Không có xét nghiệm nào cho kết quả chính xác $100\%$. Một loại xét nghiệm, mà ở đây ta gọi là xét nghiệm X, cho kết quả như sau: Khi con bò bị bệnh bò điên thì xác suất để có phản ứng dương tính trong xét nghiệm X là $70\%$, còn khi con bò không bị bệnh  thì xác suất để có phản ứng dương tính trong xét nghiệm X là $10\%$. Biết rằng tỉ lệ bò bị mắc bệnh bò điên ở Hà Lan là $13$ con trên $1\,000\,000$ con. Khi một con bò ở Hà Lan có phản ứng dương tính với xét nghiệm X thì xác suất để nó bị mắc bệnh bò điên là
	\choice 
		{$\dfrac{91}{100\,078}$} 
		{\True $\dfrac{91}{1\,000\,078}$} 
		 {$\dfrac{91}{3\,000\,052}$}
		{$\dfrac{91}{8\,999\,974}$} 
	\loigiai{ 
		 Xét các biến cố:
		$A$: \lq\lq Con bò  ở Hà Lan bị bệnh bò điên\rq\rq\;\\
		$B$: \lq\lq Con bò ở Hà Lan có phản ứng dương tính  với xét nghiệm X \rq\rq.\\
		 Theo giả thiết, ta có $\mathrm{P}(A)=0{,}000013$; $\mathrm{P}(B\mid A)=0{,}7$; $\mathrm{P}\left(B\mid \overline{A}\right)=0{,}1$.\\
		Theo công thức  xác suất toàn phần, ta có
		$$ 
		\mathrm{P}(B)=\mathrm{P}(A) \cdot \mathrm{P}(B \mid  A)+\mathrm{P}\left(\overline{A}\right) \cdot \mathrm{P}\left(B \mid \overline{A}\right)=0{,}000013 \cdot 0{,}7+(1-0{,}000013) \cdot 0{,}1=0{,}1000078.
		 $$
		Theo công thức Bayes, ta có  $\mathrm{P}(A\mid B)=\dfrac{\mathrm{P}(A) \cdot \mathrm{P}(B\mid A)}{\mathrm{P}(B)}=\dfrac{0{,}000013\cdot 0{,}7}{0{,}1000078}=\dfrac{91}{1\,000\,078}$.
	 }
\end{ex}

% nguồn: Lop_12_Dai_so_Chuong_6
\begin{ex}%[2D6H2-3]%[ma:74911c7c]
	Một bộ lọc được sử dụng để chặn thư rác trong các tài khoản thư điện tử. Tuy nhiên, vì bộ lọc không tuyệt đối hoàn hảo nên một thư rác bị chặn với xác suất $0{,}95$ và một thư đúng (không phài là thư rác) bị chặn với xác suất $0{,}01$. Thống kê cho thấy tỉ lệ thư rác là $3\%$. Chọn ngẫu  nhiên một thư bị chặn. Tính xác suất để đó là thư rác (kết quả làm tròn đến hàng phần nghìn).
	 \choice
		{$0{,}095$} 
		{\True  $0{,}746$}
		{$0{,}476$} 
		 {$0{,}003$}
	\loigiai{ 
		Gọi $A$ là biến cố \lq\lq  Thư được chọn là thư rác\rq\rq.\\
		$B$ là biến cố \lq\lq  Thư được chọn là bị chặn\rq\rq.\\
		Ta có $\mathrm{P}(A)=3\%=0{,}03$;  $\mathrm{P}\left(\overline{A}\right)=1-\mathrm{P}(A)=1-0{,}03=0{,}97$.\\
		Đồng  thời $\mathrm{P}(B \mid A)=0{,}95$; $\mathrm{P}\left(B \mid \overline{A}\right)=0{,}01$.\\
		Công thức Bayes, ta có 
		$$ 
		\mathrm{P}(A\mid B)=\dfrac{\mathrm{P}(A)  \cdot \mathrm{P}(B\mid A)}{\mathrm{P}(A) \cdot \mathrm{P}(B\mid A)+\mathrm{P}\left(\overline{A}\right) \cdot \mathrm{P}\left(B\mid \overline{A}\right)}=\dfrac{0{,}03 \cdot 0{,}95}{0{,}03 \cdot 0{,}95+0{,}97 \cdot 0{,}01} \approx 0{,}746.
		$$ 
	} 
\end{ex}

% nguồn: Lop_12_Dai_so_Chuong_6
\begin{ex}%[2D6H2-3]%[ma:d9b33400]
	Nếu hai biến  cố $A$, $B$ thỏa mãn $P(A)=0{,}4$, $P(B\mid A)=0{,}6$ thì $P(A\cap B)$ bằng
	 \choice
		{\True $\dfrac{6}{25}$} 
		 {$\dfrac{2}{3}$}
		{$\dfrac{1}{5}$} 
		{$1$} 
	\loigiai{Ta có  $P(B\mid A)=\dfrac{P(A\cap B)}{P(A)}\Rightarrow P(A\cap B)=P(B\mid A)\cdot P(A)=0{,}6\cdot 0{,}4=\dfrac{6}{25}$.}
\end{ex}

% nguồn: Lop_12_Dai_so_Chuong_6
\begin{ex}%[2D6H2-3]%[ma:f4197328]
	Cho hai biến cố $A$, $B$ thoả mãn $\mathrm{P}(A)=0{,}4$;  $\mathrm{P}(B)=0{,}3$; $\mathrm{P}(A \mid B)=0{,}25$. Khi đó, $\mathrm{P}(B \mid A)$ bằng
	 \choice
		{\True $0{,}1875$} 
		{$0{,}48$} 
		{$0{,}333$} 
		{$0{,}95$} 
	 \loigiai{
		Theo công thức Bayes, ta có $\mathrm{P}(B \mid  A)=\dfrac{\mathrm{P}(B) \cdot \mathrm{P}(A \mid B)}{\mathrm{P}(A)}=\dfrac{0{,}3 \cdot 0{,}25}{0{,}4}=0{,}1875$.
	 }
\end{ex}

% nguồn: Lop_12_Dai_so_Chuong_6
\begin{ex}%[2D6H2-3]%[ma:65c07403]
	Một ứng dụng được sử dụng để chặn cuộc gọi rác trong điện thoại. Tuy nhiên, vì ứng dụng không tuyệt đối hoàn hảo nên một cuộc gọi rác bị chặn với xác suất $0{,}8$ và một cuộc gọi đúng (không phải là cuộc gọi rác) bị chặn với xác suất $0{,}01$. Thống kê cho thấy tỉ lệ cuộc gọi rác là $10\%$. Chọn ngẫu nhiên một cuộc gọi không bị chặn. Xác suất để  đó là cuộc gọi đúng là
	 \choice
		 {$\dfrac{891}{911}$}
		{\True $\dfrac{891}{911}$} 
		{$\dfrac{123}{892}$} 
		 {$\dfrac{213}{911}$}
	\loigiai{ 
		Gọi $A$ là biến cố: \lq\lq cuộc gọi được chọn là cuộc gọi rác\rq\rq\\ $B$ là biến cố: \lq\lq  cuộc gọi chọn bị chặn\rq\rq\ thì $\overline{B}$ là biến cố: \lq\lq cuộc gọi được chọn không bị chặn\rq\rq.\\
		Theo đầu bài  ta có $\mathrm{P}(A)=0{,}1$, $\mathrm{P}\left(\overline{A}\right)=0{,}9$, $\mathrm{P}(B\mid A)=0{,}8$, $\mathrm{P}\left(B\mid \overline{A}\right)=0{,}01$.\\
		Ta có $\mathrm{P}(B)=\mathrm{P}(B\mid A) \cdot \mathrm{P}(A)+\mathrm{P}\left(B\mid \overline{A}\right) \cdot  \mathrm{P}\left(\overline{A}\right)=0{,}8\cdot 0{,}1+0{,}01\cdot 0{,}9=0{,}089$.\\
		Vì $\mathrm{P}\left(B\mid \overline{A}\right)=0{,}01\Rightarrow  \mathrm{P}\left(\overline{B} \mid \overline{A}\right)=0{,}99, \mathrm{P}(B\mid A)=0{,}8\Rightarrow \mathrm{P}\left(\overline{B} \mid A\right)=0{,}2$.\\
		Theo công  thức Bayes ta có
		 $$
		\mathrm{P}(\overline{A} \mid \overline{B})=\dfrac{\mathrm{P}\left(\overline{A}\right) \cdot \mathrm{P}\left(\overline{B} \mid  \overline{A}\right)}{\mathrm{P}\left(\overline{A}\right) \cdot \mathrm{P}\left(\overline{B} \mid \overline{A}\right)+\mathrm{P}(A) \cdot \mathrm{P}\left(\overline{B} \mid A\right)}=\dfrac{0{,}9 \cdot 0{,}99}{0{,}9 \cdot 0{,}99+0{,}1 \cdot 0{,}2}=\dfrac{891}{911}.
		 $$
	} 
\end{ex}

% nguồn: Lop_12_Dai_so_Chuong_6
\begin{ex}%[2D6H2-3]%[ma:5b04513e]
	Giả sử tỉ lệ người dân của tỉnh X nghiện thuốc lá  là $20\%$; tỉ lệ người bị bệnh phổi trong số người nghiện thuốc lá là $70\%$, trong số người không nghiện thuốc lá là $15\%$. Khi ta gặp ngẫu nhiên một người dân của tỉnh X, xác suất mà người đó là nghiện thuốc lá khi biết bị bệnh phổi là
	 \choice
		 {\True $\dfrac{7}{13}$}
		{$\dfrac{6}{13}$} 
		{$\dfrac{4}{13}$} 
		 {$\dfrac{9}{13}$}
	 \loigiai{
		Gọi $A$ là biến cố \lq\lq người đó nghiện thuốc  lá\rq\rq;\\
		 $B$ là biến cố \lq\lq người đó bị bệnh phổi\rq\rq.\\
		Ta có $\mathrm{P}(A)=0{,}2$; $\mathrm{P}\left(\overline{A}\right)=1-\mathrm{P}(A)=1-0{,}2=0{,}8$; $\mathrm{P}(B\mid A)=0{,}7$  và $\mathrm{P}\left(B\mid \overline{A}\right)=0{,}15$.\\
		Theo công thức xác suất  toàn phần ta có
		\allowdisplaybreaks 
		\begin{align*}
			 \mathrm{P}(B)&=\mathrm{P}(A)\cdot\mathrm{P}(B\mid A)+\mathrm{P}\left(\overline{A}\right)\cdot\mathrm{P}\left(B\mid\overline{A}\right)\\
			 &=0{,}2\cdot 0{,}7+0{,}8\cdot 0{,}15=0{,}26.
		\end{align*}
		Khi đó xác suất người đó nghiện thuốc lá khi biết bị bệnh phổi là $\mathrm{P}(A\mid B)$.\\ 
		Áp dụng công thức Bayes ta có $\mathrm{P}(A\mid B)=\dfrac{\mathrm{P}(A)\cdot\mathrm{P}(B\mid  A)}{\mathrm{P}(B)}=\dfrac{0{,}2\cdot 0{,}7}{0{,}26}=\dfrac{7}{13}$.
	} 
\end{ex}

% nguồn: Lop_12_Dai_so_Chuong_6
\begin{ex}%[2D6H2-3]%[ma:2f6c6580]
	Giả sử tỉ lệ người dân của thủ đô Hà Nội nghiện thuốc lá là $30\%$; tỉ lệ người bị bệnh phổi là $38\%$ và tỉ lệ người bị bệnh phổi trong số người nghiện thuốc lá là $80\%$. Chọn ngẫu nhiên một người của thủ đô Hà Nội, tính xác suất mà người đó là nghiện thuốc lá khi biết bị  bệnh phổi.
	 \choice
		 {$\dfrac{7}{13}$}
		{$\dfrac{6}{19}$} 
		 {$\dfrac{4}{13}$}
		 {\True $\dfrac{12}{19}$}
	\loigiai{ 
		Gọi $A$ là biến cố \lq\lq người  nghiện thuốc lá\rq\rq.\\
		Gọi $B$ là biến cố \lq\lq người bị  bệnh phổi\rq\rq.\\
		Ta có $\mathrm{P}(A)=0{,}3$;  $\mathrm{P}(B)=0{,}38$; $\mathrm{P}(B\mid A)=0{,}8$.\\
		Xác suất mà người đó là nghiện thuốc lá khi biết bị bệnh  phổi là $\mathrm{P}(A\mid B)$.\\
		Theo công thức Bayes, ta có $\mathrm{P}(A\mid  B)=\dfrac{\mathrm{P}(A) \cdot \mathrm{P}(B\mid A)}{\mathrm{P}(B)}=\dfrac{0{,}3\cdot 0{,}8}{0{,}38}=\dfrac{12}{19}$.\\
		Trong số người bị bệnh phổi của thủ đô Hà Nội, có khoảng $\dfrac{12}{19}$ số người  nghiện thuốc lá.
	 }
\end{ex}

% nguồn: Lop_12_Dai_so_Chuong_6
\begin{ex}%[2D6H2-3]%[ma:88fe205d]
	Kểt quả khảo sát tại một xã cho thấy có $25\%$ cư dân hút thuốc lá. Tỉ lệ cư dân thường xuyên gặp các vấn đề sức khoẻ về đường hô hấp trong số những người hút thuốc lá và không hút thuốc lá lần lượt là $60\%$ và $25\%$. Nếu ta gặp một cư dân của xã thường xuyên gặp các vấn đề sức khoẻ về đường hô hấp thì xác suất người đó có hút thuốc lá là bao nhiêu? 
	\choice 
		{\True $\dfrac{4}{9}$} 
		{$\dfrac{5}{9}$} 
		{$\dfrac{7}{9}$} 
		 {$\dfrac{8}{9}$}
	 \loigiai{
		Gọi $A$ là biến cố  \lq\lq Người đó có hút thuốc lá\rq\rq;\\
		$B$ là biến  cố \lq\lq Người đó thường xuyên gặp các vấn đề sức khoẻ về đường hô hấp\rq\rq.\\
		Ta có $\mathrm{P}(A\mid B)=\dfrac{\mathrm{P}(A) \cdot \mathrm{P}(B\mid A)}{\mathrm{P}(A)  \cdot \mathrm{P}(B\mid A)+\mathrm{P}\left(\overline{A}\right) \cdot \mathrm{P}\left(B\mid \overline{A}\right)}=\dfrac{0{,}25 \cdot 0{,}6}{0{,}25 \cdot 0{,}6+0{,}75 \cdot 0{,}25}=\dfrac{4}{9}$.\\
		Vậy nếu ta gặp một cư  dân của xã thường xuyên gặp các vấn đề sức khoẻ về đường hô hấp thì xác suất người đó có hút thuốc lá là $\dfrac{4}{9}$.
	} 
\end{ex}

% nguồn: Lop_12_Dai_so_Chuong_6
\begin{ex}%[2D6H2-3]%[ma:3ad770dc]
	Một căn bệnh X có $1\%$ dân số mắc phải. Một phương pháp chuẩn  đoán được phát triển có tỷ lệ chính xác là $99\%$. Với những người bị bệnh X, phương pháp này sẽ đưa ra kết quả dương tính $99\%$ số trường hợp. Với người không mắc bệnh, phương pháp này cũng chuẩn đoán đúng $99$ trong $100$ trường hợp. Nếu một người kiểm tra và kết quả là dương tính (bị bệnh), xác suất để người đó thực sự bị bệnh là bao nhiêu?
	\choice 
		 {$0{,}4$}
		{$0{,}35$} 
		{\True  $0{,}5$}
		 {$0{,}65$}
	 \loigiai{
		Gọi $A$ là biến cố \lq\lq người đó  mắc bệnh\rq\rq,\\
		$B$ là biến cố  \lq\lq kết quả kiểm tra người đó là dương tính (bị bệnh)\rq\rq.\\
		Xác suất  để người đó mắc bệnh khi chưa kiểm tra là $\mathrm{P}(A)=1\%=0{,}01$.\\
		Do đó xác suất đề người đó không mắc bệnh khi chưa kiểm tra là  $\mathrm{P}\left(\overline{A}\right)=1-0{,}01=0{,}99$.\\
		Xác suất kết quả dương tính nếu người  đó mắc bệnh là $\mathrm{P}(B\mid A)=99\%=0{,}99$.\\
		Xác suất kết quả dương tính nếu  người đó không mắc bệnh là $\mathrm{P}\left(B\mid \overline{A}\right)=1-0{,}99=0{,}01$.\\
		Xác  suất để người đó thực sự bị bệnh là
		$$ 
		\mathrm{P}(A\mid  B)=\dfrac{\mathrm{P}(A) \cdot \mathrm{P}(B\mid A)}{\mathrm{P}(A) \cdot \mathrm{P}(B\mid A)+\mathrm{P}\left(\overline{A}\right) \cdot \mathrm{P}\left(B\mid \overline{A}\right)}=\dfrac{0{,}01 \cdot 0{,}99}{0{,}01 \cdot 0{,}99+0{,}99 \cdot 0{,}01}=0{,}5.
		 $$
	 }
\end{ex}

% nguồn: Lop_12_Dai_so_Chuong_6
\begin{ex}%[2D6H2-3]%[ma:04455058]
	Hai máy tự động sản xuất cùng một  loại chi tiết, trong đó máy I sản xuất $35\%$, máy II sản xuất $65\%$ tổng sản lượng. Tỉ lệ phế phẩm của các máy lần lượt là $0{,}3\%$ và $0{,}7\%$. Chọn ngẫu nhiên $1$ sản phẩm từ kho. Tính xác suất để chọn được phể phẩm do máy I sản xuất.
	\choice 
		{$0{,}0056$} 
		 {\True $0{,}1875$}
		{$0{,}1785$} 
		{$0{,}1587$} 
	\loigiai{ 
		 Gọi $A_1$ là biến cố \lq\lq Sản phẩm được chọn do máy I sản xuất\rq\rq\\
		$A_2$ là biến cố \lq\lq  Sản phẩm được chọn do máy II sản xuất\rq\rq;\\
		$B$ là biến cố \lq\lq Sản phẩm  được chọn là phế phẩm\rq\rq.\\
		Suy ra $A_1\mid B$ là biến cố \lq\lq chọn được phế phẩm do máy I sản  xuất\rq\rq.\\
		Ta có $\mathrm{P}\left(A_1\right)=0{,}35$, $\mathrm{P}\left(A_2\right)=0{,}65$, $\mathrm{P}\left(B\mid A_1\right)=0{,}003$, $\mathrm{P}\left(B\mid  A_2\right)=0{,}007$.\\
		Áp dụng công thức xác suất toàn phần ta  có
		$$\mathrm{P}(B)=\mathrm{P}\left(B\mid A_1\right)  \cdot \mathrm{P}\left(A_1\right)+\mathrm{P}\left(B\mid A_2\right) \cdot \mathrm{P}\left(A_2\right)=0{,}0056.$$
		Theo công thức Bayes  có $\mathrm{P}\left(A_1\mid B\right)=\dfrac{\mathrm{P}\left(B\mid A_1\right) \cdot \mathrm{P}\left(A_1\right)}{\mathrm{P}(B)}=0{,}1875$.
	 }
\end{ex}

% nguồn: Lop_12_Dai_so_Chuong_6
\begin{ex}%[2D6H2-3]%[ma:e964df69]
	Cho hai biến cố $A$, $B$ thỏa  mãn $\mathrm{P}(A)=0{,}3$; $\mathrm{P}(B)=0{,}2$ và $\mathrm{P}(A\mid B)=0{,}15$. Khi đó, $\mathrm{P}(B\mid A)$ bằng
	 \choice
		{\True  $0{,}1$}
		 {$0{,}4$}
		{$0{,}225$} 
		 {$0{,}009$}
	 \loigiai{
		Theo  công thức Bayes, ta có $\mathrm{P}(B\mid A)=\dfrac{\mathrm{P}(B) \cdot \mathrm{P}(A\mid B)}{\mathrm{P}(A)}=\dfrac{0{,}2\cdot 0{,}15}{0{,}3}=0{,}1$.
	} 
\end{ex}

% nguồn: Lop_12_Dai_so_Chuong_6
\begin{ex}%[2D6H2-3]%[ma:519935ac]
	Cho hai biến cố  $A$ và $B$. Biết rằng $\mathrm{P}(B)=0{,}8$; $\mathrm{P}(A\mid B)=0{,}7$ và\break $\mathrm{P}\left(A\mid\overline{B}\right)=0{,}45$. Khi đó giá trị của $\mathrm{P}(B\mid A)$ bằng
	\choice 
		{$0{,}25$} 
		 {$0{,}65$}
		{\True  $\dfrac{56}{65}$}
		{$0{,}5$} 
	\loigiai{ 
		Ta có  $\mathrm{P}\left(\overline{B}\right)=1-\mathrm{P}(B)=1-0{,}8=0{,}2$.\\
		Theo công thức xác  suất toàn phần ta có
		\allowdisplaybreaks 
		\begin{align*}
			 \mathrm{P}(A)&=\mathrm{P}(B)\cdot\mathrm{P}(A\mid B)+\mathrm{P}\left(\overline{B}\right)\cdot\mathrm{P}\left(A\mid\overline{B}\right)\\
			&=0{,}8\cdot 0{,}7+0{,}2\cdot 0{,}45=0{,}65. 
		\end{align*}
		Khi đó áp  dụng công thức Bayes ta có $\mathrm{P}(B\mid A)=\dfrac{\mathrm{P}(B)\cdot\mathrm{P}(A\mid B)}{\mathrm{P}(A)}=\dfrac{0{,}8\cdot 0{,}7}{0{,}65}=\dfrac{56}{65}$.
	} 
\end{ex}

% nguồn: Lop_12_Dai_so_Chuong_6
\begin{ex}%[2D6H2-3]%[ma:4a5e9999]
	Áo sơ mi An Phước trước khi xuất khẩu sang Mỹ phải qua 2 lần  kiểm tra, nếu cả hai lần đều đạt thì chiếc áo đó mới đủ tiêu chuẩn xuất khẩu. Biết rằng bình quân $98\%$ sản phẩm làm ra qua được lần kiểm tra thứ nhất, và $95\%$ sản phẩm qua được lần kiểm tra đầu sẽ tiếp tục qua được lần kiểm tra thứ hai. Tìm xác suất để $1$ chiếc áo sơ mi đủ tiêu chuẩn xuất khẩu.
	\choice 
		{$\dfrac{95}{98}$} 
		{\True $\dfrac{931}{1000}$} 
		{$\dfrac{95}{100}$} 
		{$\dfrac{98}{100}$} 
	 \loigiai{
		Gọi $A$ là biến cố \lq\lq Qua được lần kiểm tra đầu tiên\rq\rq\,$\Rightarrow  \mathrm{P}(A)=0{,}98$.\\
		Gọi $B$ là biến cố \lq\lq Qua được lần kiểm tra thứ hai\rq\rq\,$\Rightarrow \mathrm{P}(B\mid  A)=0{,}95$.\\
		Chiếc áo sơ mi  đủ tiêu chuẩn xuất khẩu phải thỏa mãn 2 điều kiện trên hay ta đi tính $\mathrm{P}(A\cap B)$.\\
		Ta có $\mathrm{P}(A\cap B)=\mathrm{P}(B\mid A) \cdot  \mathrm{P}(A)=0{,}95\cdot 0{,}98=\dfrac{931}{1000}$.
	} 
\end{ex}

% nguồn: Lop_12_Dai_so_Chuong_6
\begin{ex}%[2D6H2-3]%[ma:5257b892]
	Giả sử có một loại bệnh S mà tỉ lệ người mắc bệnh là $0{,}1\%$. Giả sử có một loại xét nghiệm, mà ai mắc bệnh S khi xét nghiệm cũng có phản ứng dương tính, nhưng tỉ lệ phản ứng dương tính giả là $5\%$ (tức là trong số những người không bị bệnh S có $5\%$ số người xét  nghiệm lại có phản ứng dương tính). Khi một người xét nghiệm có phản ứng dương tính thì khả năng mắc bệnh S của người đó là bao nhiêu phần trăm (làm tròn kết quả đến hàng phần trăm)?
	\choice 
		{\True $1{,}96\%$} 
		{$1{,}69\%$} 
		 {$1{,}97\%$}
		{$0{,}5\%$} 
	\loigiai{ 
		Gọi $A$ là biến cố: \lq\lq Người  đó mắc bệnh S\rq\rq\;\\
		$B$ là biến cố: \lq\lq Người  đó xét nghiệm có phản ứng dương tính\rq\rq.\\
		Ta  có $\mathrm{P}(A)=0{,}001$; $\mathrm{P}\left(\overline{A}\right)=1-\mathrm{P}(A)=1-0{,}001=0{,}999$.\\
		Đồng thời  $\mathrm{P}(B\mid A)=1$; $\mathrm{P}\left(B\mid \overline{A}\right)=0{,}05$.\\
		Thay vào công  thức Bayes ta được:
		$$ 
		\mathrm{P}(A\mid B)=\dfrac{\mathrm{P}(A) \cdot \mathrm{P}(B\mid A)}{\mathrm{P}(A) \cdot \mathrm{P}(B\mid A)+\mathrm{P}\left(\overline{A}\right) \cdot \mathrm{P}\left(B\mid \overline{A}\right)}=\dfrac{0{,}001\cdot 1}{0{,}001\cdot  1+0{,}999\cdot 0{,}05}=\dfrac{20}{1019} \approx 1{,}96\%.
		$$ 
	 }
\end{ex}

% nguồn: Lop_12_Dai_so_Chuong_6
\begin{ex}%[2D6H2-3]%[ma:1d4d3579]
	Công ty $X$ tổ chức buổi trình diễn ánh sáng với $1\,000$ drone, điều khiển từ trạm trung tâm để tạo hình khối 3D đẹp mắt, sử dụng hệ thống định vị RTK (Real Time Kinematic) đảm bảo độ chính xác. Trong quá trình bay, drone có thể gặp sự cố mất tín hiệu kết nối do nhiễu sóng, và khi đó, nó tự động phát tín hiệu cảnh báo qua kênh dự phòng. Tuy nhiên, hệ thống cảnh báo không hoàn hảo, có thể báo nhầm ngay cả khi drone hoạt  động bình thường. Dựa trên thống kê từ những buổi trình diễn trước đó: Xác suất một drone gặp sự cố mất kết nối là $0{,}005$; xác suất drone bình thường bị báo nhầm phát tín hiệu cảnh báo là $0{,}001$; xác suất drone gặp sự cố tín hiệu phát tín hiệu cảnh báo là $0{,}99$. Trong buổi trình diễn, một drone được ghi nhận phát tín hiệu cảnh báo. Xác suất để drone này thực sự gặp sự cố mất tín hiệu kết nối là bao nhiêu? (làm tròn kết quả đến phần trăm).
	(Drone chuyên  trình diễn ánh sáng)
	 \choice
		 {$0{,}06$}
		{\True $0{,}83$} 
		 {$0{,}79$}
		 {$0{,}33$}
	\loigiai{Gọi  $A$ là biến cố \lq\lq drone bị mất kết nối\rq\rq, ta có $P(A)=0{,}005$ và $P\left(\overline{A}\right)=0{,}995$.\\
		Gọi $B$ là biến cố \lq\lq hệ thống cảnh nhận  được tín hiệu từ drone\rq\rq, ta có $\heva{&P(B\mid A)=0{,}99\\&P\left(B\mid\overline{A}\right)=0{,}001.}$\\
		Xác  suất cần tìm là
		\[P(A\mid B)=\dfrac{P(A)\cdot P(B\mid  A)}{P(A)\cdot P(B\mid A)+P\left(\overline{A}\right)\cdot P\left(B\mid \overline{A}\right)}=\dfrac{990}{1189}\approx0{,}83.\]
	 }
\end{ex}

% nguồn: Lop_12_Dai_so_Chuong_6
\begin{ex}%[2D6H2-3]%[ma:d6201c14]
	Cho hai biến cố $A$, $B$ thoả mãn  $\mathrm{P}(A)=0{,}4$; $\mathrm{P}(B)=0{,}3$; $\mathrm{P}(A \mid B)=0{,}25$. Khi đó, $P(B\mid A)$ bằng
	\choice 
		{$0{,}333$} 
		 {\True $0{,}1875$}
		 {$0{,}48$}
		 {$0{,}95$}
	\loigiai{Áp dụng công thức Bayes  ta được $\mathrm{P}(B \mid A)=\dfrac{\mathrm{P}(A\mid B) \cdot \mathrm{P}(B)}{\mathrm{P}(A)}=\dfrac{0{,}25 \cdot 0{,}3}{0{,}4}=0{,}1875$. }
\end{ex}


%%%%%%%%%%%%%%%% ĐÚNG/SAI %%%%%%%%%%%%%%%%

% nguồn: Lop_12_Dai_so_Chuong_6
\begin{ex}%[2D6H2-3]%[ma:7547083d]
	Trong một hộp có $8$ viên bi màu xanh và $6$ viên bi màu đỏ, các viên bi cùng kích thước và cùng khối lượng. Bạn Hùng lấy ngẫu nhiên một viên bi từ hộp, không trả lại. Sau đó bạn Nam lấy ngẫu nhiên một viên bi trong số các bi còn lại trong hộp. Gọi $A$ là biến cố:\lq\lq Hùng lấy được viên bi màu đỏ \rq\rq, $B$ là biến cố: \lq\lq Nam lấy được viên bi màu  xanh \rq\rq.
	 \choiceTF
		{Với $\Omega$  là không gian mẫu, ta có $n(\Omega)=196$}
		{\True $\mathrm{P}(AB)=\dfrac{24}{91}$} 
		{$\mathrm{P}(B)=\dfrac{8}{13}$} 
		{\True $\mathrm{P}(A \mid  B)=\dfrac{6}{13}$}
	 \loigiai{
		\begin{itemchoice}
			 \itemch Ta có $n(\Omega)=14 \cdot 13=182 \ne 196$.
			\itemch Xác suất $\mathrm{P}(AB)=\mathrm{P}(A) \cdot \mathrm{P}(B|A)=\dfrac{6}{14} \cdot \dfrac{8}{13}=\dfrac{48}{182}=\dfrac{24}{91}$. 
			\itemch 
			 Xác suất $\mathrm{P}(\overline{A}B)=\mathrm{P}(\overline{A}) \cdot \mathrm{P}(B|\overline{A})=\dfrac{4}{7} \cdot \dfrac{7}{13}=\dfrac{4}{13}$.\\
			Xác suất $\mathrm{P}(B)=\mathrm{P}(AB)+\mathrm{P}(\overline{A}B)=\dfrac{24}{91}+\dfrac{4}{13}=\dfrac{4}{7}\ne  \dfrac{8}{13}$.
			\itemch Xác suất $\mathrm{P}(A|B)=\dfrac{\mathrm{P}(AB)}{\mathrm{P}(B)}=\dfrac{\frac{24}{91}}{\frac{4}{7}}=\dfrac{24}{91} \cdot  \dfrac{7}{4}=\dfrac{6}{13}$.
		\end{itemchoice}
	} 
\end{ex}

% nguồn: Lop_12_Dai_so_Chuong_6
\begin{ex}%[2D6H2-3]%[ma:e5ddf184]
	Giả sử bệnh hiểm nghèo $X$ có tỉ lệ nhiễm bệnh là $0{,}5\%$, xét nghiệm loại bệnh này có  tỉ lệ dương tính giả là $4\%$. Khi xét nghiệm cho một người, ta gọi $A$ là biến cố \lq\lq Người được chọn không nhiễm bệnh\rq\rq\,và $B$ là biến cố \lq\lq người được chọn có phản ứng dương tính\rq\rq.
	 \choiceTF
		{\True Người được chọn không  nhiễm bệnh có tỉ lệ $\mathrm{P}(A)=0{,}995$}
		{\True Tỉ lệ người không nhiễm bệnh trong số những người có phản ứng dương tính  là $\mathrm{P}(B \mid A)=0{,}04$}
		{Tỉ lệ người nhiễm bệnh trong số những người có phản ứng dương tính là  $\mathrm{P}\left(B\mid \overline{A}\right)=0{,}005$}
		 {\True Khả năng nhiễm bệnh của một người có phản ứng dương tính là $\mathrm{P}\left(\overline{A} \mid B\right)=\dfrac{25}{224}$}
	\loigiai{ 
		\begin{itemchoice}
			\itemch  Người được chọn không mắc bệnh có tỉ lệ $\mathrm{P}(A)=1-0{,}5\%=0{,}995$.
			\itemch Do  trong số những người không mắc bệnh có $4\%$ phản ứng dương tính nên \[\mathrm{P}(B \mid A)=0{,}04.\]
			\itemch Những người mắc bệnh đều có phản ứng  dương tính nên $\mathrm{P}\left(B\mid \overline{A}\right)=1$.
			\itemch Khả năng mắc bệnh của một người có  phản ứng dương tính là
			$$\mathrm{P}\left(\overline{A} \mid B\right)=\dfrac{\mathrm{P}\left(\overline{A}\right) \mathrm{P}\left(B\mid \overline{A}\right)}{\mathrm{P}\left(\overline{A}\right) \mathrm{P}\left(B\mid \overline{A}\right)+\mathrm{P}(A) \mathrm{P}(B \mid  A)}=\dfrac{0{,}005\cdot 1}{0{,}005\cdot 1+0{,}995\cdot 0{,}04}=\dfrac{25}{224}.$$
		\end{itemchoice}
	 }
\end{ex}

% nguồn: Lop_12_Dai_so_Chuong_6
\begin{ex}%[2D6H2-3]%[ma:405afd18]
	Một hộp có $4$ viên bi xanh và $6$ viên bi đỏ, các viên bi có kích thước và khối lượng như nhau. Lấy lần lượt  hai viên bi, không hoàn lại.
	\choiceTF 
		{Xác  suất lần $1$ lấy được bi xanh là $\dfrac{1}{4}$}
		{Xác suất lần $2$ lấy được  bi xanh biết lần $1$ lấy được bi đỏ là $\dfrac{1}{3}$}
		{Xác suất lần $2$ lấy được bi xanh biết lần $1$ lấy được  bi xanh là $\dfrac{4}{9}$}
		{\True Xác suất lần  $2$ lấy được bi xanh là $\dfrac{2}{5}$}
	 \loigiai{Gọi $A$ là biến cố "Lần $1$ lấy được bi xanh".\\
		$B$ là biến cố "Lần $2$ lấy được bi xanh, lần $1$ lấy được  bi đỏ".\\
		$C$  là biến cố "Lần $2$ lấy được bi xanh, lần $1$ lấy được bi xanh".\\
		$D$ là biến cố "Lần  $2$ lấy được bi xanh".\\
		\begin{itemchoice}
			\itemch \textbf{Sai}. Xác suất lần  $1$ lấy được bi xanh là $\mathrm{P}(A)=\dfrac{4}{10}=\dfrac{2}{5}$.
			\itemch \textbf{Sai}. Sau lần $1$ trong hộp còn  $9$ viên trong đó có có $4$ viên bi xanh và $5$ viên bi đỏ. Nên xác suất lần $2$ lấy được bi xanh biết lần $1$ lấy được bi đỏ là $\mathrm{P}(B)=\dfrac{4}{9}$.
			\itemch \textbf{Sai}. Sau lần $1$ trong hộp còn $9$ viên trong  đó có có $3$ viên bi xanh và $6$ viên bi đỏ. Nên xác suất lần $2$ lấy được bi xanh biết lần $1$ lấy được bi xanh là $\mathrm{P}(C)=\dfrac{3}{9}=\dfrac{1}{3}$.
			\itemch \textbf{Đúng}.  Ta có
			\begin{center}
				$\mathrm{P}(D)=\mathrm{P}(A)\cdot \mathrm{P}(C)+\mathrm{P}(\overline{A})\cdot \mathrm{P}(B)=\dfrac{4}{10}\cdot  \dfrac{3}{9}+\dfrac{6}{10}\cdot\dfrac{4}{9}=\dfrac{2}{5}$.
			\end{center}
			Xác  suất lần $2$ lấy được bi xanh là $\dfrac{2}{5}$.
		\end{itemchoice}
	 }
\end{ex}

% nguồn: Lop_12_Dai_so_Chuong_6
\begin{ex}%[2D6H2-3]%[ma:b6d73588]
	Một căn bệnh  có $2\%$ dân số mắc phải. Một phương pháp chẩn đoán được phát triển có tỷ lệ chính xác là $99\%$. Với những người bị bệnh, phương pháp này sẽ đưa ra kết quả dương tính $99\%$ sổ trường hợp. Với người không mắc bệnh, phương pháp này cũng chuẩn đoán đúng $97\%$. Lấy một người đi kiểm tra.
	\choiceTF 
		{\True Xác suất để  người đó mắc bệnh khi chưa kiểm tra là $0{,}02$}
		{\True Xác suất kết  quả dương tính nếu người đó mắc bệnh là $0{,}99$}
		{Xác suất kết quả dương tính nếu  người đó không mắc bệnh là $0{,}01$}
		{Biết rằng đã có kết quả chuẩn đoán là  dương tính, xác suất để người đó thực sự bị bệnh là $0{,}25$}
	\loigiai{ 
		\begin{itemchoice}
			\itemch Gọi $A$ là biến cố \lq\lq người đó mắc bệnh\rq\rq.\\ 
			Xác suất để người đó mắc bệnh khi chưa  kiểm tra $\mathrm{P}(A)=2\%=0{,}02$.
			 \itemch Gọi $B$ là biến cố \lq\lq kết quả kiểm tra người đó là dương tính\rq\rq\\
			Xác suất kết quả  dương tính nếu người đó mắc bệnh là
			 \[\mathrm{P}(B\mid A)=99\%=0{,}99.\]
			\itemch Xác xuất  kết quả âm tính nếu người đó không mắc bệnh là $\mathrm{P}\left(\overline{B} \mid\overline{A}\right)=97\%=0{,}97$.\\
			Xác suất kết quả dương tính nếu người đó không mắc bệnh  là
			 \[\mathrm{P}\left(B\mid \overline{A}\right)=1-\mathrm{P}\left(\overline{B} \mid \overline{A}\right)=1-0{,}97=0{,}03.\]
			\itemch Do đó xác suất để người đó không mắc bệnh khi chưa kiểm  tra $\mathrm{P}\left(\overline{A}\right)=1-0{,}02=0{,}98$.\\
			Xác suất để người đó thực  sự bị bệnh là
			\[\mathrm{P}(A \mid B)=\dfrac{\mathrm{P}(A) \cdot \mathrm{P}(B \mid  A)}{\mathrm{P}(A) \cdot \mathrm{P}(B \mid A)+\mathrm{P}\left(\overline{A}\right) \cdot \mathrm{P}\left(B\mid \overline{A}\right)}=\dfrac{0{,}02 \cdot 0{,}99}{0{,}02 \cdot 0{,}99+0{,}98 \cdot 0{,}03}=\dfrac{33}{82} \approx 0{,}402.\]
		\end{itemchoice}
	 }
\end{ex}

% nguồn: Lop_12_Dai_so_Chuong_6
\begin{ex}%[2D6H2-3]%[ma:6bdcd0f7]
	Chạy Marathon là môn thể thao mà tại đó, người chơi sẽ hoàn thành quãng đường $42{,}195$\,km trong khoảng thời gian nhất định. FM sub $4$ là thành tích dành cho những người chơi hoàn thành quãng đường Marathon dưới $4$ giờ. Trong CLB AKR, tỷ lệ thành viên nam là $72\%$, tỷ lệ thành viên nữ là $28\%$. Đối với nam, tỷ lệ Vận động viên hoàn thành Marathon sub $4$ là $32\%$; đối với nữ tỷ  lệ Vận động viên hoàn thành sub $4$ là $3\%$. Chọn ngẫu nhiên $1$ thành viên từ CLB AKR.
	 \choiceTF
		{\True Khi Vận động viên được chọn là nam, xác suất để Vận động viên này chưa hoàn thành sub $4$ cự ly Marathon  là $68\%$}
		{Xác suất để  thành viên được chọn đã hoàn thành sub $4$ là $22\%$}
		{Xác suất để thành viên được chọn là  nữ đã hoàn thành sub $4$ là $2\%$}
		{\True Biết  rằng Vận động viên được chọn đã hoàn thành sub $4$, xác suất để Vận động viên đó là nam bằng $96\%$}
	\loigiai{ 
		Gọi $A$ là biến cố \lq\lq  Vận động viên được chọn là nam\rq\rq.\\
		Gọi $B$ là biến cố \lq\lq Vận động viên được chọn đã  hoàn thành cự ly Marathon sub $4$\rq\rq.
		\begin{itemchoice}
			\itemch sub $4$ Khi vận động  viên được chọn là nam, xác suất để vận động viên này chưa hoàn thành sub $4$ cự ly Marathon là
			$$\mathrm{P}\left(\overline{B} \mid A\right)=1-\mathrm{P}(B  \mid A)=1-32\%=68\%.$$
			\itemch Xác suất  để vận động viên được chọn đã hoàn thành sub $4$ là
			$$\mathrm{P}(B)=\mathrm{P}(A) \cdot \mathrm{P}(B \mid A)+\mathrm{P}\left(\overline{A}\right) \cdot  \mathrm{P}\left(B\mid \overline{A}\right)=0{,}72\cdot 0{,}32+0{,}28\cdot 0{,}03 \approx 0{,}24=24 \%. $$
			\itemch Xác suất để vận động viên được chọn là  nữ và đã hoàn thành sub $4$ là
			 $$\mathrm{P}\left(\overline{A}B\right)=\mathrm{P}\left(\overline{A}\right) \cdot \mathrm{P}\left(B\mid \overline{A}\right)=0{,}28 \cdot 0{,}03 \approx 0{,}0084 \approx 0{,}84 \%.$$
			\itemch Biết vận động viên đã  hoàn thành sub $4$, xác suất để vận động viên đó là nam là
			 \allowdisplaybreaks
			\begin{align*}
				\mathrm{P}(A \mid  B)&=\dfrac{\mathrm{P}(A) \cdot \mathrm{P}(B \mid A)}{\mathrm{P}(B)}\\
				&=\dfrac{\mathrm{P}(A) \cdot \mathrm{P}(B \mid A)}{\mathrm{P}(A) \cdot \mathrm{P}(B \mid  A)+\mathrm{P}\left(\overline{A}\right) \cdot \mathrm{P}\left(B\mid \overline{A}\right)} \\
				 &=\dfrac{0{,}72 \cdot 0{,}32}{0{,}72 \cdot 0{,}32+0{,}28 \cdot 0{,}03} \approx 0{,}96=96\%.
			\end{align*}
		\end{itemchoice}
	 }
\end{ex}

% nguồn: Lop_12_Dai_so_Chuong_6
\begin{ex}%[2D6H2-3]%[ma:add73d17]
	Một lớp có $70\%$ học sinh là nữ.  Tỉ lệ học sinh nữ đạt danh hiệu học sinh giỏi là $35\%$, tỉ lệ học sinh nam đạt danh hiệu học sinh giỏi là $60\%$. Chọn ngẫu nhiên một học sinh của lớp đó. Gọi $A$ là biến cố \lq \lq Học sinh được chọn là nữ\rq \rq và $B$ là biến cố \lq \lq Học sinh được chọn đạt danh hiệu học sinh giỏi\rq \rq.
	\choiceTF 
		{Xác suất  của biến cố $\overline{A}$ là $0{,}7$}
		{\True Xác suất  của biến cố $B$ là $0{,}425$}
		{$A$ và  $B$ là hai biến cố độc lập}
		{Xác suất của biến  cố $A$ với điều kiện $B$ bằng $\dfrac{5}{7}$}
	\loigiai{ 
		\begin{itemchoice}
			\itemch Vì  $\mathrm{P}(A)=0{,}7$ nên $\mathrm{P}(\overline{A})=0{,}3$.
			 \itemch Ta có $\mathrm{P}(B|A)=0{,}35$ và $\mathrm{P}(B|\overline{A})=0{,}6$.\\
			Khi đó $\mathrm{P}(B)=\mathrm{P}(A)\cdot \mathrm{P}(B|A)+\mathrm{P}(\overline{A})\cdot \mathrm{P}(B|\overline{A})=0{,7}\cdot  0{,}35+0{,}3\cdot 0,6=0{,}425$.
			\itemch Để  $A$ và $B$ độc lập thì $\mathrm{P}(B|A)=\mathrm{P}(B)$.\\
			 Vì $\mathrm{P}(B|A) \ne \mathrm{P}(B)$ nên $A$ và $B$ không độc lập.
			 \itemch $\mathrm{P}(A|B)=\dfrac{\mathrm{P}(B|A)\cdot \mathrm{P}(A)}{\mathrm{P}(B)}=\dfrac{0{,}35 \cdot 0{,}7}{0{,}425}=\dfrac{49}{85}$.
		\end{itemchoice}
	 }
\end{ex}

% nguồn: Lop_12_Dai_so_Chuong_6
\begin{ex}%[2D6H2-3]%[ma:2e9cb170]
	Một doanh nghiệp có $45\%$ nhân viên là nữ. Tỉ lệ nhân viên nữ và tỉ lệ nhân viên nam mua bảo hiểm nhân thọ lần lượt là $7\%$ và $5\%$. Chọn  ngẫu nhiên một nhân viên của doanh nghiệp.
	\choiceTF 
		{Xác suất nhân viên được chọn có mua  bảo hiểm nhân thọ là $0{,}061$}
		{\True Biết rằng nhân viên được chọn có mua  bảo hiểm nhân thọ. Xác suất nhân viên đó là nam là $\dfrac{55}{118}$}
		{\True Biết rằng nhân viên được chọn có mua bảo hiểm nhân thọ. Xác suất nhân viên đó là  nữ là $\dfrac{63}{118}$}
		{Biết rẳng nhân viên được chọn có  mua bảo hiểm nhân thọ. Khi đó nhân viên đó là nam nhiều hơn là nữ}
	\loigiai{ 
		Gọi $A$ là biến cố \lq\lq Nhân viên được chọn là nữ\rq\rq\ và $B$ là biến cố \lq\lq Nhân viên được  chọn có mua bảo hiểm nhân thọ\rq\rq.\\
		Theo đề ta có $\mathrm{P}(A)=0{,}45$; $\mathrm{P}(B\mid A)=0{,}07$; $\mathrm{P}\left(B\mid  \overline{A}\right)=0{,}05$.\\
		Suy ra $\mathrm{P}\left(\overline{A}\right)=0{,}55$. 
		\begin{itemchoice}
			\itemch Ta có $\mathrm{P}(B)=\mathrm{P}(A) \cdot \mathrm{P}(B\mid A)+\mathrm{P}\left(\overline{A}\right) \cdot \mathrm{P}\left(B\mid \overline{A}\right)=0{,}45\cdot  0{,}07+0{,}55\cdot 0{,}05=0{,}059$.
			\itemch $\mathrm{P}(\overline{A} \mid B)=\dfrac{\mathrm{P}\left(\overline{A}\right) \cdot \mathrm{P}\left(B\mid \overline{A}\right)}{\mathrm{P}(B)}=\dfrac{0{,}55\cdot 0{,}05}{0{,}059}=\dfrac{55}{118}$. 
			\itemch $\mathrm{P}(A\mid B)=\dfrac{\mathrm{P}(A) \cdot \mathrm{P}(B\mid A)}{\mathrm{P}(B)}=\dfrac{0{,}45\cdot 0{,}07}{0{,}059}=\dfrac{63}{118}$. 
			\itemch $\mathrm{P}(A\mid B)=\dfrac{\mathrm{P}(A)  \cdot \mathrm{P}(B\mid A)}{\mathrm{P}(B)}=\dfrac{0{,}45.0{,}07}{0{,}059}=\dfrac{63}{118}$.\\
			Do $\mathrm{P}(A\mid B)=\dfrac{63}{118} > \dfrac{55}{118}=\mathrm{P}\left(\overline{A} \mid B\right)$ nên nhân viên được chọn có mua bảo hiểm nhân thọ là nữ sẽ nhiều hơn là nam. 
		\end{itemchoice}
	} 
\end{ex}

% nguồn: Lop_12_Dai_so_Chuong_6
\begin{ex}%[2D6H2-3]%[ma:5241e372]
	Một chiếc hộp có $50$  viên bi, trong đó có $30$ viên bi màu đỏ và $20$ viên bi màu vàng; các viên bi có kích thước và khối lượng như nhau. Sau khi kiểm tra, người ta thấy có $80\%$ số viên bi màu đỏ đánh số và $60\%$ số viên bi màu vàng có đánh số, những viên bi còn lại không đánh số. Lấy ra ngẫu nhiên một viên bi trong hộp.
	\choiceTF 
		{\True Xác suất để lấy được bi đánh số có màu  vàng là $0{,}6$}
		{Xác suất để  lấy được bi không đánh số có màu đỏ là $0{,}8$}
		{Xác suất để viên bi  được lấy ra có đánh số là $0{,}36$}
		{\True Xác suất  đề lấy viên bi màu đỏ có đánh số là $\dfrac{2}{3}$}
	 \loigiai{
		Gọi $A$  là biến cố \lq\lq viên bi được lấy ra có đánh số\rq\rq\\
		Gọi $B$ là biến cố \lq\lq viên bi được lấy ra có màu đỏ\rq\rq\, suy ra $\overline{B}$ là biến cố \lq\lq viên  bi được lấy ra có màu vàng\rq\rq\,
		\begin{itemchoice}
			\itemch $\mathrm{P}\left(A\mid \overline{B}\right)=60\%=0{,}6$. 
			 \itemch $\mathrm{P}(A\mid B)=80\%=0{,}8$, nên $\mathrm{P}\left(\overline{A} \mid B\right)=1-\mathrm{P}(A\mid B)=1-0{,}8=0{,}2$.
			\itemch  Ta có $\mathrm{P}(B)=\dfrac{30}{50}=0{,}6$; $\mathrm{P}\left(\overline{B}\right)=\dfrac{20}{50}=0{,}4$.\\
			Vậy $\mathrm{P}(A)=\mathrm{P}(B) \cdot \mathrm{P}(A\mid B)+\mathrm{P}\left(\overline{B}\right)  \cdot \mathrm{P}\left(A\mid \overline{B}\right)=0{,}6\cdot 0{,}8+0{,}4\cdot 0{,}6=0{,}72$.
			 \itemch Ta có $\mathrm{P}(B\mid A)=\dfrac{\mathrm{P}(B) \mathrm{P}(A\mid B)}{\mathrm{P}(A)}=\dfrac{0{,}6\cdot 0{,}8}{0{,}72}=\dfrac{2}{3}$.
		\end{itemchoice}
	} 
\end{ex}

% nguồn: Lop_12_Dai_so_Chuong_6
\begin{ex}%[2D6H2-3]%[ma:111f0276]
	Lớp 12A có $70\%$ học sinh thích chơi thể thao. Biết rằng, nếu một học sinh thích chơi thể thao thì xác  xuất học sinh đó biết chơi cầu lông là $0{,}8$; học sinh không thích chơi thể thao thì xác suất học sinh đó biết chơi cầu lông là $0{,}1$. Chọn ngẫu nhiên một học sinh.
	 \choiceTF
		{\True Xác suất học sinh  này không thích chơi thể thao là $0{,}3$}
		{Xác suất học sinh này  không biết chơi cầu lông với điều kiện không thích chơi thể thao là $0{,}41$}
		{\True Xác suất  học sinh này biết chơi cầu lông là $0{,}59$}
		{\True Xác suất học sinh này thích chơi thể thao với điều kiện biết chơi cầu lông là $0{,}95$ (làm tròn kết  quả đến hàng phần trăm)}
	\loigiai{Gọi $A$ là biến  cố \lq\lq Học sinh thích chơi thể thao\rq\rq, ta có $P(A)=0{,}7$.\\
		Gọi $B$ là biến cố \lq\lq Học  sinh biết chơi cầu lông\rq\rq, ta có $\heva{&P(B\mid A)=0{,}8\\&P\left(B\mid\overline{A}\right)=0{,}1.}$
		\begin{itemchoice}
			\itemch Xác suất học sinh được chọn  không thích chơi thể thao là $P\left(\overline{A}\right)=0{,}3$.
			\itemch Xác suất học sinh  này không biết chơi cầu lông với điều kiện không thích chơi thể thao là
			 \[P\left(\overline{B}\mid\overline{A}\right)=1-P\left(B\mid\overline{A}\right)=1-0{,}1=0{,}9.\]
			\itemch Xác suất học  sinh này biết chơi cầu lông là
			\[P(B)=P(A)\cdot P(B\mid A)+P\left(\overline{A}\right)\cdot  P\left(B\mid\overline{A}\right)=0{,}59.\]
			\itemch  Xác suất học sinh này thích chơi thể thao với điều kiện biết chơi cầu lông là
			\[P(A\mid B)=\dfrac{P(A)\cdot P(B\mid A)}{P(A)\cdot P(B\mid A)+P\left(\overline{A}\right)\cdot  P\left(B\mid\overline{A}\right)}=\dfrac{56}{59}\approx0{,}95.\]
	\end{itemchoice}}
\end{ex}

% nguồn: Lop_12_Dai_so_Chuong_6
\begin{ex}%[2D6H2-3]%[ma:e134638d]
	Trong cuộc khảo sát ngẫu nhiên $600$ học sinh về việc \lq\lq thích học\rq\rq\ hay \lq\lq không thích học\rq\rq\ môn toán. Kết quả thống kê như sau: có $400$ học sinh trả lời \lq\lq thích học\rq\rq\ và $200$ học sinh trả lời \lq\lq không thích học\rq\rq\. Thực tế cho thấy tỉ lệ học sinh  \lq\lq giỏi toán\rq\rq\ tương ứng với những cách trả lời \lq\lq thích học\rq\rq\ và \lq\lq không thích học\rq\rq\ lần lượt là $70\%$ và $10\%$.
	Gọi $T$ là biến cố \lq\lq Học sinh được phỏng vấn trả lời thích học môn toán\rq\rq. 
	Gọi $G$ là biến cố \lq\lq Học sinh được phỏng  vấn giỏi toán\rq\rq.
	\choiceTF 
		{\True Xác suất $\mathrm{P}(T) = \dfrac{2}{3}$ và $\mathrm{P}(\overline{T})  = \dfrac{1}{3}$}
		{Xác suất có điều kiện $\mathrm{P}(G|T) = 0{,}3$,  $\mathrm{P}(G|\overline{T}) = 0{,}9$}
		{\True Xác suất $\mathrm{P}(G)  = 0{,}5$}
		{Trong số những người được phỏng vấn giỏi toán có $80\%$ người đã trả lời thích học môn toán khi được phỏng vấn (kết quả tính theo phần trăm làm tròn  đến hàng đơn vị)}
	 \loigiai{
		Gọi  $T$ là biến cố \lq\lq Học sinh được phỏng vấn trả lời thích học môn toán\rq\rq.\\
		\noindent Gọi  $G$ là biến cố \lq\lq Học sinh được phỏng vấn giỏi toán\rq\rq\.
		Theo giả thiết $\mathrm{P}(T) = \dfrac{2}{3}$, $\mathrm{P}(\overline{T}) = \dfrac{1}{3}$, $\mathrm{P}(G|T) = 0{,}7$  và $\mathrm{P}(G|\overline{T}) = 0{,}1$.
		\begin{itemchoice}
			\itemch Ta  có $\mathrm{P}(T) = \dfrac{2}{3},\ \mathrm{P}(\overline{T}) = \dfrac{1}{3}$.
			\itemch  Ta có $\mathrm{P}(G|T) = 0{,}7$ và $\mathrm{P}(G|\overline{T}) = 0{,}1$.
			\itemch Ta có $$ 
			\mathrm{P}(G)=\mathrm{P}(T)\cdot\mathrm{P}(G|T)+\mathrm{P}(\overline{T})\cdot\mathrm{P}(G|\overline{T})= \dfrac{2}{3}\cdot  0{,}7+\dfrac{1}{3}\cdot 0{,}1=\dfrac{1}{2}.$$
			 \itemch Theo công thức Bayes
			$$\mathrm{P}(T|G)=\dfrac{\mathrm{P}(T\cap G)}{\mathrm{P}(G)}  =\dfrac{\mathrm{P}(T)\cdot\mathrm{P}(G|T)}{\mathrm{P}(T)\cdot\mathrm{P}(G|T)+\mathrm{P}(\overline{T})\cdot\mathrm{P}(G|\overline{T})}
			=  \dfrac{0{,}7 \cdot \dfrac{2}{3}}{\dfrac{2}{3}\cdot 0{,}7+\dfrac{1}{3}\cdot 0{,}1} = \dfrac{14}{15}\neq 0{,}8.$$
		\end{itemchoice}
	} 
\end{ex}

% nguồn: Lop_12_Dai_so_Chuong_6
\begin{ex}%[2D6V2-3]%[ma:aefbfe81]
	Nghiên cứu số bệnh nhân trong một viện bỏng, thấy rằng có $2$ nguyên nhân gây ra bỏng là bỏng nhiệt và bỏng do hóa chất. Bỏng nhiệt chiếm $70\%$ số bệnh nhân và bỏng do hóa chất là $30\%$. Trong những bệnh nhân bị bỏng nhiệt thì có $30\%$ bị biến chứng, trong những bệnh nhân bị bỏng  hóa chất thì có $50\%$ bị biến chứng. Rút ngẫu nhiên một bệnh án.
	\choiceTF 
		{\True Xác  suất của bỏng nhiệt bị biến chứng là $0{,}3$}
		{\True Xác suất của bỏng hóa chất bị biến chứng  là $0{,}5$}
		{Xác suất của bệnh án bị  biến chứng là $32\%$}
		{\True Biết rằng bệnh  án rút ra bị biến chứng, xác suất bệnh án đó do bỏng nhiệt là $\dfrac{7}{12}$}
	 \loigiai{
		Gọi $A$ là biến cố \lq\lq bệnh án rút ra bị biến  chứng\rq\rq.\\
		Gọi $B$ là biến  cố \lq\lq bệnh án rút ra bị bỏng nhiệt\rq\rq.\\
		Khi  đó $\overline{B}$ là biến cố \lq\lq bệnh án rút ra bị bỏng hóa chất\rq\rq.
		\begin{itemchoice}
			 \itemch
			Xác suất  do bị bỏng nhiệt là $\mathrm{P}(B)=0{,}7$.\\
			Xác suất bị biến chứng trong bỏng nhiệt  là $30\%=0{,}3\Rightarrow \mathrm{P}(A\mid B)=0{,}3$.\\
			Xác suất do bị bỏng  hóa chất là $\mathrm{P}(\overline{B})=0{,}3$.
			\itemch 
			Xác suất bị biến chứng trong bỏng hóa chất là $50\%=0{,}5\Rightarrow \mathrm{P}(A\mid\overline{B})=0{,}5$. 
			\itemch Ta  có
			\begin{align*}
				\mathrm{P}(A)=\mathrm{P}(AB)+\mathrm{P}(A\overline{B})&=\mathrm{P}(B)\cdot\mathrm{P}(A\mid B)+\mathrm{P}(\overline{B})\cdot \mathrm{P}(A\mid\overline{B})\\ 
				&=0{,}3\cdot 0{,}7+0{,}5\cdot  0{,}3=36\%.
			\end{align*}
			 \itemch
			Xác suất của bệnh án bị biến chứng do  bỏng nhiệt
			\[\mathrm{P}( B\mid  A)=\dfrac{\mathrm{P}(B)\cdot\mathrm{P}(A\mid B)}{\mathrm{P}(A)}=\dfrac{0{,}7\cdot0{,}3}{0{,}36}=\dfrac{7}{12}.\]
		\end{itemchoice}
	} 
\end{ex}

% nguồn: Lop_12_Dai_so_Chuong_6
\begin{ex}%[2D6V2-3]%[ma:350d3c2a]
	Có hai đội thi đấu môn bắn súng. Đội I có $8$ vận động viên, đội II có $10$ vận động viên. Xác suất đạt huy chương vàng của mỗi vận động viên đội I và đội II tương ứng là  $0{,}6$ và $0{,}55$. Chọn ngẫu nhiên một vận động viên.
	\choiceTF 
		{Xác suất để vận động  viên chọn ra thuộc đội I là $\dfrac{5}{9}$}
		{\True Xác suất không đạt huy chương vàng của mỗi  vận động viên đội II là $0{,}45$}
		{\True Xác suất để vận động viên này đạt huy chương vàng  là $\dfrac{103}{180}$}
		 {\True Giả sử vận động viên được chọn đạt huy chương vàng. Xác suất để vận động viên này thuộc đội I là $\dfrac{48}{103}$}
	 \loigiai{
		\begin{itemchoice}
			\itemch Xác  suất để vận động viên chọn ra thuộc đội I là $\dfrac{8}{18}=\dfrac{4}{9}$.
			\itemch Xác suất không  đạt huy chương vàng của mỗi vận động viên đội II là
			 \[1-0{,}55=0{,}45.\]
			\itemch Gọi $A$ là biến cố \lq\lq Vận động viên  đạt huy chương vàng\rq\rq\\
			$B$ là biến cố \lq\lq Thành viên đội I\rq\rq\ thì biến cố đối  của $B$ là $\overline{B}\colon$ \lq\lq Thành viên đội II đạt huy chương vàng\rq\rq.\\
			 Do đó, $\mathrm{P}(B)=\dfrac{8}{18}=\dfrac{4}{9}$; $\mathrm{P}\left(\overline{B}\right)=\dfrac{5}{9}$; $\mathrm{P}(A\mid B)=0{,}6; \mathrm{P}\left(A\mid \overline{B}\right)=0{,}55$.\\
			 Theo công thức xác suất toàn phần ta có
			\[\mathrm{P}(A)=\mathrm{P}(B) \cdot \mathrm{P}(A \mid B)+\mathrm{P}\left(\overline{B}\right) \cdot \mathrm{P}(A \mid  \overline{B})=\dfrac{4}{9} \cdot 0{,}6+\dfrac{5}{9} \cdot 0{,}55=\dfrac{103}{180}.\]
			\itemch Ta có  $\mathrm{P}(B\mid A)=\dfrac{\mathrm{P}(B) \cdot \mathrm{P}(A\mid B)}{\mathrm{P}(A)}=\dfrac{\dfrac{4}{9} \cdot 0{,}6}{\dfrac{103}{180}}=\dfrac{48}{103}$.
		\end{itemchoice}
	} 
\end{ex}

% nguồn: Lop_12_Dai_so_Chuong_6
\begin{ex}%[2D6V2-3]%[ma:bf4a018b]
	Một công ty công nghệ cung  cấp hai phiên bản Basic và Pro của một phần mềm. Tỉ lệ người sử dụng hai phiên bản này lần lượt là $80\%$ và $20\%$. Kết quả điều tra cho thấy có $30\%$ người dùng
	phiên bản  Basic sẽ mua bản cập nhật sau $1$ năm sử dụng; còn tỉ lệ này của phiên bản Pro là $60\%$.
	Chọn ngẫu nhiên  một người sử dụng phần mềm trên của công ty.
	Gọi $A$ là biến cố: \lq\lq Người dùng mua bản cập nhật sau $1$  năm sử dụng\rq\rq\,và $B$ là biến cố: \lq\lq Người dùng
	sử dụng phiên bản  Basic ở năm đầu tiên\rq\rq.
	\choiceTF 
		 {\True $\mathrm{P}(A|B)=0{,}3 $}
		{$\mathrm{P}(AB)=0{,}12$} 
		{$\mathrm{P}(A)=0{,}24  $}
		{\True $\mathrm{P}(B|A)=\dfrac{2}{3}  $}
	\loigiai{ 
		\begin{itemchoice}
			\itemch  Đúng.\\Vìcó $30\%$ người dùng phiên bản Basic sẽ mua bản cập nhật sau $1$ năm sử dụng
			 nên $\mathrm{P}(A|B)=0{,}3$.
			\itemch Sai.\\ 
			Theo công thức nhân xác suất  ta có
			 $$\mathrm{P}(AB)=\mathrm{P}(B)\cdot \mathrm{P}(A|B)=0{,}8\cdot 0{,}3=0{,}24.$$
			 \itemch Sai.\\
			 Áp dụng công thức xác suất toàn phần, ta có
			$$\mathrm{P}(A)=\mathrm{P}(B)\cdot  \mathrm{P}(A|B)+\mathrm{P}(\overline{B})\cdot \mathrm{P}(A|\overline{B})=0{,}8\cdot 0{,3}+0{,}2\cdot 0{,}6=0{,}36.$$
			\itemch  Đúng.\\
			Áp dụng công  thức Bayes, ta có $$\mathrm{P}(B|A)=\dfrac{\mathrm{P}(B)\cdot \mathrm{P}(A|B)}{\mathrm{P}(B)\cdot \mathrm{P}(A|B)+\mathrm{P}(\overline{B})\cdot \mathrm{P}(A|\overline{B})}=\dfrac{0{,}24}{0{,}36}=\dfrac{2}{3}.$$
		\end{itemchoice}
	} 
\end{ex}

% nguồn: Lop_12_Dai_so_Chuong_6
\begin{ex}%[2D6V2-3]%[ma:99e769b4]
	Một loại xét nghiệm nhậnh SARS-CoV-2 cho kết quả dương tính với $76{,}2\%$ các ca thực sự nhiễm virus và kết quả âm tính với $99{,}1\%$ các ca  thực sự không nhiễm virus. Giả sử tỉ lệ người nhiễm virus SARS-CoV-2 trong một cộng đồng là $1\%$.
	\choiceTF 
		{Xác suất xét nghiệm cho kết quả âm tính của các ca thực  sự nhiễm virus là $0{,}23$}
		{\True Xác suất xét nghiệm cho kết quả dương tính của các ca thực sự không nhiễm virus là $0{,}009$} 
		{\True Xác suất người làm  xét nghiệm có kết quả dương tính là $0{,}017$}
		{\True Biết rằng đã có kết quả chuẩn đoán là dương tính, xác suất để người  đó thực sự bị bệnh là $\dfrac{381}{850}$}
	 \loigiai{
		Gọi $A$ là biến cố \lq\lq Người làm xét nghiệm  có kết quả dương tính\rq\rq\ và $B$ là biến cố \lq\lq Người nhiễm virus\rq\rq.
		\begin{itemchoice}
			\itemch Do xét nghiệm  cho kết quả dương tính với $76{,}2\%$ các ca thực sự nhiễm virus nên xác suất xét nghiệm cho kết quả âm tính của các ca thực sự nhiễm virus là $\mathrm{P}(\overline{A} \mid B)=0{,}238$.
			\itemch Do xét nghiệm cho kết quả âm tính với $99{,}1\%$ các ca thực sự không nhiễm virus  nên $\mathrm{P}(\overline{A} \mid \overline{B})=0{,}991$.\\
			Suy ra $\mathrm{P}\left(A\mid  \overline{B}\right)=1-0{,}991=0{,}009$.\\
			Xác suất xét nghiệm cho kết quả dương tính của các ca thực  sự không nhiễm virus là $0{,}009$.
			\itemch Do tỉ lệ người nhiễm virus trong cộng đồng  là $1\%$ nên $\mathrm{P}(B)=0{,}01$ và\\ $\mathrm{P}\left(\overline{B}\right)=0{,}99$.\\
			Áp dụng công thức xác  suất toàn phần, ta có xác suất người làm xét nghiệm có kết quả dương tính là
			 \[\mathrm{P}(A)=\mathrm{P}(B) \mathrm{P}(A\mid B)+P\left(\overline{B}\right) \mathrm{P}\left(A\mid \overline{B}\right)=0{,}01\cdot 0{,}762+0{,}99\cdot 0{,}009=0{,}01653\approx 0{,}017.\]
			\itemch Xác suất để người đó thực sự bị bệnh khi có kết quả chuẩn đoán là  dương tính là
			\[\mathrm{P}(B\mid A)=\dfrac{\mathrm{P}(B)  \mathrm{P}(A\mid B)}{\mathrm{P}(B) \mathrm{P}(A\mid B)+\mathrm{P}\left(\overline{B}\right) \mathrm{P}\left(A\mid \overline{B}\right)}=\dfrac{\mathrm{P}(B) \mathrm{P}(A\mid B)}{\mathrm{P}(A)}=\dfrac{0{,}01\cdot 0{,}762}{0{,}017}=\dfrac{381}{850}.\]
		\end{itemchoice}
	} 
\end{ex}

% nguồn: Lop_12_Dai_so_Chuong_6
\begin{ex}%[2D6V2-3]%[ma:3a0f6284]
	Trong một kì thi tốt nghiệp trung học phổ thông, một tỉnh $X$ có $50\%$ học sinh lựa chọn tổ hợp B00 (Gồm các môn Toán, Hóa, Sinh). Biết rằng, nếu một học sinh chọn tổ hợp B00 thì xác suất để học sinh đó đỗ đại học là $0{,}6$; còn nếu một học sinh không chọn tổ hợp B00 thì xác suất để học sinh đó đỗ đại học là $0{,}7$. Chọn ngẫu nhiên một học sinh của tỉnh $X$ đã tốt  nghiệp trung học phổ thông trong kì thi trên. Biết rằng học sinh này đã đỗ đại học. Gọi $A$ là biến cố: \lq\lq Học sinh đó chọn tổ hợp B00\rq\rq\; $B$ là biến cố: \lq\lq Học sinh đó đỗ đại học\rq\rq.
	\choiceTF 
		{\True Xác  suất $\mathrm{P}\left(\overline{A}\right)=0{,}5$}
		 {Xác suất $\mathrm{P}(B\mid A)=0{,}4$}
		{\True Xác suất $\mathrm{P}\left(B\mid \overline{A}\right)$ thuộc khoảng  $(0{,}2; 0{,}5)$}
		 {$\dfrac{\mathrm{P}(A\mid B)}{\mathrm{P}(B\mid A)}$ lớn hơn $\dfrac{2}{3}$}
	 \loigiai{
		\begin{itemchoice}
			\itemch  $\overline{A}$ là biến cố: \lq\lq Học sinh đó không chọn tổ hợp B00\rq\rq.\\
			Ta có $\mathrm{P}(A)=0{,}5\Rightarrow \mathrm{P}\left(\overline{A}\right)=1-\mathrm{P}(A)=1-0{,}5=0{,}5$. 
			\itemch $\mathrm{P}(B\mid A)$ là xác suất để một học sinh đỗ đại học với điều kiện học  sinh đó chọn tổ hợp B00 $\Rightarrow \mathrm{P}(B\mid A)=0{,}6$.
			\itemch $\mathrm{P}\left(B\mid \overline{A}\right)$ là xác suất để một học sinh đỗ đại học với điều kiện học sinh đó không chọn tổ hợp B00.  Ta có $\mathrm{P}\left(B\mid \overline{A}\right)=0{,}7$.
			\itemch Theo công thức Bayes, ta có $\mathrm{P}(A\mid B)=\dfrac{\mathrm{P}(A) \cdot \mathrm{P}(B\mid A)}{\mathrm{P}(A) \mathrm{P}(B\mid A)+\mathrm{P}\left(\overline{A}\right) \mathrm{P}\left(B\mid \overline{A}\right)}$.\\ 
			Suy ra $\dfrac{\mathrm{P}(A\mid B)}{\mathrm{P}(B\mid A)}=\dfrac{\mathrm{P}(A)}{\mathrm{P}(A) \mathrm{P}(B\mid  A)+\mathrm{P}\left(\overline{A}\right) \mathrm{P}\left(B\mid \overline{A}\right)}=\dfrac{0{,}5}{0{,}5\cdot 0{,}6+0{,}5\cdot 0{,}7}=\dfrac{10}{13}$.
		\end{itemchoice}
	} 
\end{ex}

% nguồn: Lop_12_Dai_so_Chuong_6
\begin{ex}%[2D6V2-3]%[ma:ea82bf9c]
	Hộp thứ  nhất có $1$ viên bi xanh và $5$ viên bi đỏ. Hộp thứ hai có $3$ viên bi xanh và $5$ viên bi đỏ. Các viên bi là khác nhau. Lấy ra ngẫu nhiên đồng thời $2$ viên bi từ hộp thứ nhất chuyển sang hộp thứ hai. Sau đó lại lấy ra ngẫu nhiên $2$ viên bi từ hộp thứ hai.
	 \choiceTF
		{\True Xác suất để hai viên bi lấy ra từ hộp  hai là bi đỏ bằng $\dfrac{19}{45}$}
		{Xác suất để hai viên  bi lấy ra từ hộp hai có $1$ bi đỏ và $1$ bi xanh bằng $\dfrac{1}{9}$}
		{\True Biết  rằng hai viên bi lấy ra từ hộp thứ hai là bi đỏ. Xác suất để $2$ viên bi lấy ra từ hộp thứ nhất cũng là bi đỏ bằng $\dfrac{14}{19}$}
		{Biết rằng hai viên bi lấy ra từ hộp thứ hai có $1$  bi đỏ và $1$ bi xanh. Xác suất để $2$ viên bi lấy ra từ hộp thứ nhất cũng có $1$ bi đỏ và $1$ bi xanh bằng $\dfrac{5}{19}$}
	 \loigiai{
		Gọi $A$ là biến cố \lq\lq Lấy được hai viên bi đỏ từ hộp thứ nhất\rq\rq\,và $B$ là biến cố \lq\lq Lấy được hai viên bi đỏ từ hộp thứ  hai\rq\rq.\\
		Gọi $C$ là biến cố \lq\lq Lấy được  $1$ viên bi đỏ và $1$ viên bi xanh từ hộp thứ nhất \rq\rq\,và $D$ là biến cố \lq\lq Lấy được $1$ viên bi đỏ và $1$ viên bi xanh từ hộp thứ hai\rq\rq.\\
		 Ta có $\mathrm{P}(A)=\dfrac{\mathrm{C}_5^2}{\mathrm{C}_6^2}=\dfrac{2}{3}$; $\mathrm{P}(B\mid A)=\dfrac{\mathrm{C}_7^2}{\mathrm{C}_{10}^2}=\dfrac{7}{15}$;\\
		Đồng  thời $\mathrm{P}(C)=\dfrac{\mathrm{C}_1^1\mathrm{C}_5^1}{\mathrm{C}_6^2}=\dfrac{1}{3}$; $\mathrm{P}(D\mid C)=\dfrac{\mathrm{C}_4^1\mathrm{C}_6^1}{\mathrm{C}_{10}^2}=\dfrac{8}{15}$.\\
		Suy ra $\mathrm{P}\left(\overline{A}\right)=1-\mathrm{P}(A)=\dfrac{1}{3}$; $\mathrm{P}\left(B\mid  \overline{A}\right)=\dfrac{\mathrm{C}_6^2}{\mathrm{C}_{10}^2}=\dfrac{1}{3}$;\\
		Ngoài ra $\mathrm{P}\left(\overline{C}\right)=1-\mathrm{P}(C)=\dfrac{2}{3}$; $\mathrm{P}\left(D\mid \overline{C}\right)=\dfrac{\mathrm{C}_3^1\mathrm{C}_7^1}{\mathrm{C}_{10}^2}=\dfrac{7}{15}$. 
		\begin{itemchoice}
			\itemch $\mathrm{P}(B)=\mathrm{P}(A) \cdot \mathrm{P}(B\mid  A)+\mathrm{P}\left(\overline{A}\right) \cdot \mathrm{P}\left(B\mid \overline{A}\right)=\dfrac{2}{3} \cdot \dfrac{7}{15}+\dfrac{1}{3} \cdot \dfrac{1}{3}=\dfrac{19}{45}$.
			\itemch $\mathrm{P}(D)=\mathrm{P}(C)  \cdot \mathrm{P}(D\mid C)+\mathrm{P}\left(\overline{C}\right) \cdot \mathrm{P}\left(D\mid \overline{C}\right)=\dfrac{1}{3} \cdot \dfrac{8}{15}+\dfrac{2}{3} \cdot \dfrac{7}{15}=\dfrac{22}{45}$.
			\itemch Áp dụng công thức Bayes ta  có $\mathrm{P}(A\mid B)=\dfrac{\mathrm{P}(A) \cdot \mathrm{P}(B\mid A)}{\mathrm{P}(B)}=\dfrac{2}{3} \cdot \dfrac{7}{15}: \dfrac{19}{45}=\dfrac{14}{19}$.
			\itemch $\mathrm{P}(C\mid D)=\dfrac{\mathrm{P}(C)  \cdot \mathrm{P}(D\mid C)}{\mathrm{P}(D)}=\dfrac{1}{3} \cdot \dfrac{8}{15}: \dfrac{22}{45}=\dfrac{4}{11}$.
		\end{itemchoice}
	} 
\end{ex}


%%%%%%%%%%%%%%%% TRẢ LỜI NGẮN %%%%%%%%%%%%%%%%

% nguồn: Lop_12_Dai_so_Chuong_6
\begin{bt}%[2D6H2-3]%[ma:5c3cab2a]
	Một nhà máy có hai phân xưởng I và II tương ứng làm ra $ 40\%$ và $60\%$ sản phẩm của nhà máy. Biết rằng tỷ lệ phế phẩm của hai phân xưởng I và II tương ứng là $1\%$  và $2\%$. Chọn ngẫu nhiên một sản phẩm của nhà máy thì thấy nó là phế phẩm. Tính xác suất để sản phẩm đó thuộc phân xưởng I. \shortans{0,25}
	\loigiai{ 
		Gọi  $A$ là biến cố \lq\lq sản phẩm lấy ra là phế phẩm\rq\rq.\\
		Gọi $B$ là biến cố  \lq\lq sản phẩm lấy ra thuộc phân xưởng I\rq\rq.\\
		Theo dữ kiện bài toán, xác suất để sản phẩm lấy ra thuộc phân xưởng I biết rằng sản phẩm  đó là phế phẩm, ta cần tính $\mathrm{P}(B\mid A)$.\\
		Theo công thức Bayes thì $\mathrm{P}(B\mid A)=\dfrac{\mathrm{P}(B)\cdot \mathrm{P}(A\mid  B)}{\mathrm{P}(A)}$.\\
		Ta có \begin{itemize}
		\item  $\mathrm{P}(B)=40\%=0{,}4$;
		\item $\mathrm{P}(A\mid  B)=1\%=0{,}01$;
		 \item $\mathrm{P}(A)=\mathrm{P}(B)\cdot \mathrm{P}(A\mid B)+\mathrm{P}\left(\overline{B}\right)\cdot \mathrm{P}\left(A\mid\overline{B}\right)=0{,}4\cdot 0{,}01+0{,}6\cdot 0{,}02=0{,}016$.
	\end{itemize}
	Vậy $\mathrm{P}(B\mid A)=\dfrac{\mathrm{P}(B)\cdot \mathrm{P}(A\mid B)}{\mathrm{P}(A)}=\dfrac{0{,}4\cdot 0{,}01}{0{,}016}=\dfrac{1}{4}=0{,}25 $. 
} 
\end{bt}

% nguồn: Lop_12_Dai_so_Chuong_6
\begin{bt}%[2D6H2-3]%[ma:aad414df]
	Một loại linh kiện do hai nhà máy I và II cùng sản xuất. Tỉ lệ phế phẩm của các nhà máy I và II lần lượt là 2\% và 3\%. Trong một lô linh kiện để lẫn lộn $100$ sản phẩm của nhà máy I và $150$ sản phẩm của nhà máy II. Một nhân viên kiểm tra lấy ngẫu nhiên một linh kiện từ lô hàng đó. Biết rằng linh kiện được lấy ra không là phế phẩm. Tính xác  suất để linh kiện đó do nhà máy II sản xuất (kết quả làm tròn đến hàng phần mười)? \shortans{0,6}
	\loigiai{ 
		Gọi $B_1$ là biến  cố: \lq \lq Linh kiện lấy ra do nhà máy I sản xuất\rq \rq.\\
		$B_2$ là biến cố: \lq \lq Linh kiện lấy ra do nhà máy II sản xuất\rq  \rq.\\
		$A$ là biến cố: \lq \lq Linh kiện  được lấy ra không phải là phế phẩm\rq \rq.\\
		Theo bài ra, ta có $\mathrm{P}(B_1)=\dfrac{100}{250}=0{,}4$;  \quad $\mathrm{P}(B_2)=\dfrac{150}{250}=0{,}6$; \quad $\mathrm{P}(A|B_1)=0{,}98$; \quad $\mathrm{P}(A|B_2)=0{,}97$.\\
		Áp dụng công thức xác  suất Bayes ta có
		$$\mathrm{P}(B_2|A)=\dfrac{\mathrm{P}(B_2).\mathrm{P}(A|B_2)}{\mathrm{P}(B_1).\mathrm{P}(A|B_1)+\mathrm{P}(B_2).\mathrm{P}(A|B_2)}=\dfrac{0{,}6\cdot 0{,}97}{0{,}4\cdot  0{,}98+0{,}6\cdot 0{,}97}=0{,}598=0{,}6.$$
	} 
\end{bt}

% nguồn: Lop_12_Dai_so_Chuong_6
\begin{bt}%[2D6H2-3]%[ma:564b9ccf]
	Một nhà máy sản xuất sản phẩm $A$ có tỷ lệ sản phẩm bị lỗi là $2\%$. Nhà máy sử dụng hai hệ thống kiểm tra chất lượng độc lập để phát  hiện lỗi:
	\begin{itemize}
		\item Hệ thống 1: Xác suất phát hiện chính xác sản phẩm lỗi là $95\%$. Xác  suất báo lỗi nhầm trên một sản phẫm không lỗi là $1\%$.
		 \item Hệ thống 2: Xác suất phát hiện chính xác sản phẩm lỗi là $90\%$. Xác suất báo lỗi nhầm trên một sản phẩm không lỗi là $5\%$.
	\end{itemize}
	Chọn ngẫu nhiên một sản phẩm. Biết rằng sản phẩm này bị cả hai hệ thống kiểm tra đều báo lỗi. Tính xác suất để sản phẩm  này thực tế không bị lỗi. Kết quả xác suất này sau khi đã làm tròn đến hàng phần nghìn là số có dạng $0{,}0ab$. Tính giá trị của $a+b$. \shortans{10}
	 \loigiai{
		Gọi $A$  là biến cố chọn được sản phẩm lỗi.\\
		Gọi $M$, $N$ lần  lượt là biến cố hệ thống $1$ kết luận sản phẩm lỗi, hệ thống $2$ kết luận sản phẩm lỗi.\\
		Ta cần tính xác suất có điều kiện:  $\mathrm{P}\left(\overline{A}\mid MN\right)=\dfrac{\mathrm{P}\left(\overline{A}MN\right)}{\mathrm{P}\left(MN\right)}$.\\
		Ta có $\mathrm{P}\left(A\right)=0{,}02\Rightarrow \mathrm{P}\left(\overline{A}\right)-0{,}98$ và $\mathrm{P}\left(M\mid A\right)=0{,}95$; $\mathrm{P}\left(M\mid\overline{A}\right)=0{,}01$. 
		$\mathrm{P}\left(N\mid AM\right)=0{,}9$;  $\mathrm{P}\left(N\mid\overline{A}M\right)=0{,}05$.\\
		 Ta có $\mathrm{P}\left(MN\right)=\mathrm{P}\left(AMN\cup \overline{A}MN\right)=\mathrm{P}\left(AMN\right)+\mathrm{P}\left(\overline{A}MN\right)$.\\
		Mặt khác $$\mathrm{P}\left(\overline{A}MN\right)=\mathrm{P}\left(\overline{A}\right)\cdot\mathrm{P}\left(M\mid\overline{A}\right)\cdot\mathrm{P}\left(N\mid\overline{A}M\right)=0{,}98\cdot 0{,}01\cdot 0{,}05.$$ và $$\mathrm{P}\left(AMN\right)=\mathrm{P}\left(A\right)\cdot  \mathrm{P}\left(M\mid A\right)\cdot\mathrm{P}\left(N\mid AM\right)=0{,}02\cdot0{,}95\cdot0{,}9.$$
		 Vậy $$\mathrm{P}\left(\overline{A}\mid MN\right)=\dfrac{0{,}98\cdot 0{,}01\cdot 0{,}05}{0{,}98\cdot 0{,}01\cdot 0{,}05+0{,}02\cdot0{,}95\cdot0{,}9}=\dfrac{49}{1759}\approx 0{,}028.$$
		Khi đó $ a+b=10$. 
	 }
\end{bt}

% nguồn: Lop_12_Dai_so_Chuong_6
\begin{bt}%[2D6H2-3]%[ma:50522807]
	Tỉ lệ bị bệnh cúm tại một địa phương bằng $0{,}25$. Khi thực hiện xét nghiệm chẩn đoán, nếu người có bệnh cúm thì khả năng  phản ứng dương tính là $96\%$, nếu người không bị bệnh cúm thì khả năng phản ứng dương tính $8\%$. Chọn ngẫu nhiên 1 người tại địa phương đó. Xác suất người được chọn có phản ứng dương tính là bao nhiêu.\shortans{0,3}
	 \loigiai{
		 Gọi $A$: \lq\lq Người được chọn bị cúm \rq\rq; $B$: \lq\lq Người được chọn bị dương tính \rq\rq.\\
		Theo đề ta có $\mathrm{P}\left(A\right)=0{,}25;\quad\mathrm{P}\left(\overline{A}\right)=0{,}75;\quad\mathrm{P}\left(b\mid  A\right)=0{,}96;\quad\mathrm{P}\left(B\mid\overline{A}\right)=0{,}08$.\\
		Ta có $$\mathrm{P}\left(B\right)=\mathrm{P}\left(A\right)\cdot\mathrm{P}\left(B\mid A\right)+\mathrm{P}\left(\overline{A}\right)\cdot\mathrm{P}\left(B\mid\overline{A}\right)=0{,}25\cdot0{,}96+0{,}75\cdot0{,}08=0{,}3.$$ 
	} 
\end{bt}

% nguồn: Lop_12_Dai_so_Chuong_6
\begin{bt}%[2D6H2-3]%[ma:bf4cbec2]
	Thực hiện  khảo sát tại một địa phương mà số trẻ em nam gấp $1{,}5$ lần số trẻ em nữ, có $8\%$ số trẻ em nam bị hen phế quản, $5\%$ số trẻ em nữ bị hen phế quảng. Chọn ngẫu nhiên một trẻ em. Giả sử trẻ em được chọn bị hen phế quản. Xác suất chọn được trẻ em nam là bao nhiêu ? (làm tròn kết quả đến hàng phần mười).\shortans{0,7}
	\loigiai{Gọi $T$: \lq\lq Trẻ em được chọn là nam \rq\rq và $H$: \lq\lq Trẻ em  được chọn bị hen phế quản \rq\rq.\\
		Theo đề bài ta có $\mathrm{P}(T)=\dfrac{3}{5};\quad\mathrm{P}(\overline{T})=\dfrac{2}{5};\quad\mathrm{P}(H\mid  T)=0{,}08;\mathrm{P}(H\mid\overline{T})=0{,}05$.\\
		Theo  công thức Bayes, ta có
		$$\mathrm{P}(T\mid H)=\dfrac{\mathrm{P}(H\mid T)\cdot  \mathrm{P}(T)}{\mathrm{P}(H\mid T)\cdot\mathrm{P}(T)+\mathrm{P}(H\mid \overline{T})\cdot\mathrm{P}(\overline{T})}=\dfrac{0{,}08\cdot\dfrac{3}{5}}{0{,}08\cdot\dfrac{3}{5}+0{,}05\cdot\dfrac{2}{5}}=\dfrac{12}{17}\approx
		0{,}7.$$ 
	 }
\end{bt}

% nguồn: Lop_12_Dai_so_Chuong_6
\begin{bt}%[2D6H2-3]%[ma:1c0d3369]
	Một căn bệnh có $1\%$ dân số mắc phải. Một phương pháp chẩn đoán căn bệnh nói trên có tỉ lệ chính xác là $98\%$ (với cả người bị bệnh và người không bị bệnh). Biết rằng nếu một người được sử dụng phương pháp trên để kiểm tra và cho kết quả dương tính (bị bệnh) thì xác suất  người đó thực sự bị bệnh là $\dfrac{y}{296}$, $y$ là số tự nhiên. Hỏi $y$ bằng bao nhiêu? \shortans{98}
	\loigiai{ 
		Gọi $A$: \lq\lq Người đó mắc bệnh \rq\rq và $B$: \lq\lq Người đó có  kết quả kiểm tra dương tính \rq\rq.\\
		Theo đề  bài ta có $\mathrm{P}\left(A\right)=0{,}01;\quad\mathrm{P}\left(B\mid A\right)=0{,}98;\quad \mathrm{P}\left(\overline{A}\right)=0{,}99;\quad\mathrm{P}\left(B\mid \overline{A}\right)=0{,}02$.\\
		Ta  có
		$$\mathrm{P}\left(A\mid B\right)=\dfrac{\mathrm{P}\left(A\right)\cdot\mathrm{P}\left(B\mid A\right)}{\mathrm{P}\left(A\right)\cdot\mathrm{P}\left(B\mid A\right)+\mathrm{P}\left(\overline{A}\right)\cdot\mathrm{P}\left(B\mid  \overline{A}\right)}=\dfrac{0{,}01\cdot 0{,}98}{0{,}01\cdot 0{,}98+0{,}99\cdot 0{,}02}=\dfrac{49}{148}.$$
		Suy ra $y=98$. 
	 }
\end{bt}

% nguồn: Lop_12_Dai_so_Chuong_6
\begin{bt}%[2D6H2-3]%[ma:33668df7]
	Tỷ lệ người nghiện thuốc lá tại một vùng là $30\%$. Biết rằng tỷ lệ người bị viêm họng trong số những người nghiện thuốc là $60\%$, còn tỷ lệ người bị viêm họng trong số những người không nghiện là $40\%$. Lấy ngẫu nhiên một người thấy người ấy không bị viêm họng. Tính xác suất người đó  nghiện thuốc lá (làm tròn kết quả tới hàng phần trăm)
	\shortans{0,22} 
	 \loigiai{
		Gọi $A$ là biến cố  \lq\lq Người này bị nghiện thuốc lá\rq\rq.\\
		$B$ là biến cố  \lq\lq Người này không bị viêm họng\rq\rq.\\
		Ta có  $\mathrm{P}(A) = 0{,}3 \Rightarrow \mathrm{P}\big(\overline{A}\big) = 0{,}7$.\\
		Tỷ lệ người bị viêm họng trong số người bị nghiện thuốc là $\mathrm{P}\big(\overline{B}  \mid A\big) = 0{,}6$.\\
		 Tỷ lệ người bị viêm họng trong số người không nghiện thuốc là $\mathrm{P}\big(\overline{B} \mid \overline{A}\big) = 0{,}4$.\\
		 Do đó
		\[ 
		\mathrm{P}\big(\overline{B}\big) 
		= \mathrm{P}(A) \cdot \mathrm{P}\big(\overline{B} \mid A\big) 
		+ \mathrm{P}\big(\overline{A}\big) \cdot \mathrm{P}\big(\overline{B} \mid  \overline{A}\big)
		= 0{,}3 \cdot 0{,}6 + 0{,}7 \cdot 0{,}4 = 0{,}46. 
		\] 
		Suy ra $\mathrm{P}(B) = 1 - \mathrm{P}\big(\overline{B}\big) =  0{,}54$.\\
		 Áp dụng công thức Bayes, ta có
		 \[
		\mathrm{P}\big(A  \mid B\big)
		 = \dfrac{\mathrm{P}(A) \cdot \mathrm{P}\big(B \mid A\big)}{\mathrm{P}(B)}
		=  \dfrac{0{,}3 \cdot 0{,}4}{0{,}54}
		= \dfrac{2}{9}  \approx 0{,}22.
		 \]
	 }
\end{bt}

% nguồn: Lop_12_Dai_so_Chuong_6
\begin{bt}%[2D6H2-3]%[ma:18f7e263]
	Một nhà máy lắp ráp nhận được các chi tiết do hai máy sản xuất. Trung bình máy thứ nhất cung cấp $65\%$ chi tiết, máy thứ hai cung cấp $35\%$ chi tiết. Khoảng $80\%$  chi tiết do máy thứ nhất sản xuất là đạt tiêu chuẩn, còn $85\%$ chi tiết do máy thứ hai sản xuất là đạt tiêu chuẩn. Lấy ngẫu nhiên từ nhà máy một sản phẩm, thấy nó đạt tiêu chuẩn. Tìm xác suất để sản phẩm đó do máy thứ nhất sản xuất (kết quả làm tròn đến hàng phần trăm).
	\shortans{0,64} 
	\loigiai{ 
		Gọi $A$  là biến cố \lq\lq Chi tiết lấy từ dây chuyền đạt tiêu chuẩn\rq\rq.\\
		$B_1$ là biến cố \lq\lq Chi tiết do  máy thứ nhất sản xuất\rq\rq.\\
		$B_2$ là biến cố \lq\lq  Chi tiết do máy thứ hai sản xuất\rq\rq.\\
		Ta  cần tính xác suất $\mathrm{P}(B_1 \mid A)$.\\
		Theo công thức Bayes,  ta có
		\[\mathrm{P}(B_1 \mid A)=\dfrac{\mathrm{P}(B_1)\cdot\mathrm{P}(A\mid B_1)}{\mathrm{P}(B_1)\cdot\mathrm{P}(A\mid  B_1)+\mathrm{P}(B_2)\cdot\mathrm{P}(A\mid B_2)}.\]
		Theo điều kiện bài toán,  ta có
		\[\mathrm{P}(B_1) = 0{,}65;\, \mathrm{P}(B_2) = 0{,}35;\, \mathrm{P}(A\mid B_1)  = 0{,}8;\, \mathrm{P}(A\mid B_2) = 0{,}85.\]
		Vậy $\mathrm{P}(B_1 \mid  A)=\dfrac{0{,}65\cdot 0{,}8}{0{,}65\cdot 0{,}8+0{,}35\cdot 0{,}85} \approx 0{,}64$.
	} 
\end{bt}

% nguồn: Lop_12_Dai_so_Chuong_6
\begin{bt}%[2D6H2-3]%[ma:dd0be8ba]
	Chạy Marathon là môn thể thao chạy bộ đường dài mà tại đó, người chơi sẽ hoàn thành  quãng đường $42{,}195$ km trong khoảng thời gian nhất định.
	\lq\lq FM sub $4$\rq\rq\ là một thuật ngữ phổ biến trong cộng đồng những người tham gia chạy Marathon, nó dùng để chỉ thành tích hoàn thành quãng đường $42{,}195$ km dưới $4$ giờ. Trong một câu lạc bộ Marathon, tỉ lệ thành viên nam là $72\%$, tỉ lệ thành viên nữ là $28\%$. Đối với nam, tỉ lệ người hoàn thành FM sub $4$ là $32\%$; đối với nữ, tỉ lệ người hoàn thành FM sub $4$ là $3\%$. Chọn ngẫu nhiên một người từ câu lạc bộ đó. Xác suất để người được chọn là nam bằng  bao nhiêu (kết quả được làm tròn đến hàng phần trăm), biết rằng người được chọn đã hoàn thành FM sub $4$? \shortans{0,96}
	\loigiai{ 
		Gọi $A$ là biến cố \lq\lq chọn  được thành viên nam\rq\rq.\\
		 Gọi $B$ là biến cố \lq\lq chọn được thành viên đã hoàn thành FM sub $4$\rq\rq.\\
		Ta có $\mathrm{P}(A)=0{,}72$; $\mathrm{P}\left(\overline{A}\right)=0{,}28$; $\mathrm{P}(B\mid A)=0{,}32$;  $\mathrm{P}\left(B\mid \overline{A}\right)=0{,}03$.\\
		Xác suất để chọn được người nam trong  câu lạc bộ biết người đó đã hoàn thành FM sub $4$ là
		\[\mathrm{P}(A\mid  B)=\dfrac{\mathrm{P}(A)\cdot \mathrm{P}(B\mid A)}{\mathrm{P}(A)\cdot \mathrm{P}(B\mid A)+\mathrm{P}\left(\overline{A}\right)\cdot \mathrm{P}\left(B\mid \overline{A}\right)}=\dfrac{0{,}32\cdot 0{,}72}{0{,}32\cdot 0{,}72+0{,}28\cdot 0{,}03}=\dfrac{192}{199}\approx 0{,}96.\]
	 }
\end{bt}

% nguồn: Lop_12_Dai_so_Chuong_6
\begin{bt}%[2D6H2-3]%[ma:161186bb]
	Trong một đợt kiểm tra sức khỏe tại trường có $200$ học sinh được xét nghiệm một loại virus.  Trong đó, biết rằng có $80$ bạn thật sự bị nhiễm virus. Nếu một bạn bị nhiễm, thì xét nghiệm cho kết quả dương tính (tức là phát hiện đúng bệnh) với xác suất $90\%$. Nếu một bạn không bị nhiễm, thì xét nghiệm vẫn có thể báo nhầm là dương tính (gọi là dương tính giả), với xác suất $5\%$. Giả sử một bạn có kết quả xét nghiệm dương tính. Hỏi xác suất để bạn đó thật sự bị nhiễm virus là bao nhiêu (làm tròn kết quả đến hàng phần trăm)? \shortans{0,92}
	 \loigiai{
		Gọi $N$ là biến cố \lq\lq Học sinh được chọn bị nhiễm virus\rq\rq.\\ 
		Gọi $D$ là  biến cố \lq\lq Học sinh được chọn có kết quả xét nghiệm dương tính\rq\rq.\\
		Ta có \begin{itemize}
		\item $\mathrm{P}(N)=\dfrac{80}{200}\Rightarrow  \mathrm{P}\left(\overline{N}\right)=1-\dfrac{80}{200}=\dfrac{120}{200}$;
		\item 
		$\mathrm{P}(D\mid N)=0{,}9$; 
		\item $\mathrm{P}\left(D\mid \overline{N}\right)=0{,}05$; 
		 \item $\mathrm{P}(D)=\dfrac{80}{200}\cdot 0{,}9+\dfrac{120}{200}\cdot 0{,}05=0{,}39$.
	\end{itemize}
	 Vậy
	$\mathrm{P}(N\mid D)=\dfrac{\mathrm{P}(N)\cdot \mathrm{P}(D\mid N)}{\mathrm{P}(D)}=\dfrac{\dfrac{80}{200}\cdot  0{,}9}{0{,}39}=0{,}92$.
 }
\end{bt}

% nguồn: Lop_12_Dai_so_Chuong_6
\begin{bt}%[2D6H2-3]%[ma:3c5cfbc3]
	Có hai chiếc hộp, hộp I có $5$ quả bóng màu  trắng và $7$ quả bóng màu đỏ, hộp II có $10$ quả bóng màu trắng và $15$ quả bóng màu đỏ, các quả bóng có cùng kích thước và khối lượng.
	Lấy ngẫu nhiên hai quả bóng từ hộp I bỏ vào hộp II. Sau đó, lấy ra ngẫu nhiên một quả bóng từ hộp II. Xác suất để quả bóng được lấy ra từ hộp II là quả bóng được chuyển từ hộp I sang, biết rằng quả bóng đó  có màu trắng là $\dfrac{a}{b}$ (là phân số tối giản). Tính $a+b$. \shortans{14}
	\loigiai{ 
		Gọi biến cố $A$: \lq\lq Hai quả bóng lấy  ra từ hộp I khác màu\rq\rq.\\
		Gọi biến cố $B$: \lq\lq Quả bóng màu trắng lấy ra từ  hộp II\rq\rq.\\
		Khi đó \begin{itemize}
		\item $\mathrm{P}(A)=\dfrac{5\cdot 7}{\mathrm{C}_{12}^{2}}=\dfrac{35}{66}$; 
		 \item $\mathrm{P}\left(\overline{A}\right)=\dfrac{\mathrm{C}_{5}^{2}}{\mathrm{C}_{12}^{2}}+\dfrac{\mathrm{C}_{7}^{2}}{\mathrm{C}_{12}^{2}}=\dfrac{5}{33}+\dfrac{7}{22}=\dfrac{31}{66}$;
		 \item $\mathrm{P}(B)=\mathrm{P}(B\mid A)\cdot \mathrm{P}(A)+\mathrm{P}\left(B\mid \overline{A}\right)\cdot \mathrm{P}\left(\overline{A}\right)=\dfrac{35}{66}\cdot\dfrac{11}{27}+\dfrac{5}{33}\cdot\dfrac{12}{27}+\dfrac{7}{22}\cdot\dfrac{10}{27}=\dfrac{65}{162}$.
	\end{itemize}
	Xác suất để quả bóng màu trắng lấy ra từ hộp II là quả bóng được chuyển  từ hộp I là
	\[\mathrm{P}(A\mid  B)=\dfrac{\mathrm{P}(B\mid A)\cdot \mathrm{P}(A)}{\mathrm{P}(B)}=\dfrac{\dfrac{1}{27}\cdot\dfrac{35}{66}+\dfrac{2}{27}\cdot\dfrac{5}{33}}{\mathrm{P}(B)}=\dfrac{1}{13}.\]
	Vậy $a=1$; $b=13$  và $a+b=14$.
} 
\end{bt}

% nguồn: Lop_12_Dai_so_Chuong_6
\begin{bt}%[2D6H2-3]%[ma:47e82ab6]
	Một xét nghiệm Covid-19 cho kết quả dương tính với $90\%$ các trường hợp thực sự nhiễm virus và cho kết quả âm tính với $80\%$ các trường hợp thực sự không nhiễm virus. Biết rằng tỉ lệ người  nhiễm Covid-19 trong một cộng đồng nào đó là $1\%$. Một người trong cộng đồng đó cho kết quả xét nghiệm dương tính. Xác suất để người đó thực sự bị nhiễm virus có dạng $\dfrac{a}{b}$ (phân số tối giản). Giá trị của $a+b$ bằng bao nhiêu?
	\shortans{24} 
	\loigiai{ 
		Gọi $A$ là biến cố \lq\lq Người  đó bị nhiễm virus\rq\rq.\\
		Gọi $B$ là biến cố \lq\lq Người đó cho kết quả dương tính\rq\rq.\\ 
		Xét nghiệm Covid-19 cho kết quả dương tính với $90\%$ các trường hợp thực  sự nhiễm virus nên $\mathrm{P}\big(B \mid A\big) = 0{,}9$.\\
		Xét nghiệm Covid-19 cho kết quả âm tính với  $80\%$ các trường hợp thực sự không nhiễm virus, nên cho kết quả dương tính với $20\%$ các trường hợp không thực sự nhiễm virus $\mathrm{P}\big(B \mid \overline{A}\big) = 0{,}2$.\\
		Ta có $\mathrm{P}(A) = 0{,}01  \Rightarrow \mathrm{P}\big(\overline{A}\big) = 0{,}99$.\\
		 Do đó, xác suất để người đó cho kết quả dương tính là
		 \[
		 \mathrm{P}(B)
		= \mathrm{P}(A) \cdot \mathrm{P}\big(B  \mid A\big)
		+ \mathrm{P}\big(\overline{A}\big) \cdot \mathrm{P}\big(B  \mid \overline{A}\big)
		= 0{,}01 \cdot 0{,}9 + 0{,}99 \cdot  0{,}2 = 0{,}207.
		\] 
		Xác suất để  người đó thực sự bị nhiễm virus khi cho kết quả dương tính là
		\[ 
		\mathrm{P}\big(A \mid B\big) 
		= \dfrac{\mathrm{P}(A) \cdot \mathrm{P}\big(B \mid  A\big)}{\mathrm{P}(B)}
		= \dfrac{0{,}01 \cdot  0{,}9}{0{,}207}
		= \dfrac{1}{23}. 
		 \]
		Vậy $a = 1$, $b = 23 \Rightarrow a + b = 24$. 
	 }
\end{bt}

% nguồn: Lop_12_Dai_so_Chuong_6
\begin{bt}%[2D6H2-3]%[ma:2c3b26a5]
	 Dây chuyền lắp ráp nhận được các chi tiết do hai máy sản xuất. Trung bình máy thứ nhất cung cấp $60\%$ chi tiết, còn lại là của máy thứ hai. Khoảng $90\%$ chi tiết do máy thứ nhất sản xuất đạt tiêu chuẩn, còn $85\%$ chi tiết do máy thứ hai sản xuất là đạt tiêu chuẩn. Lấy ngẫu nhiên từ dây chuyền một sản phẩm, thấy nó đạt tiêu chuẩn. Tính xác suất để sản phẩm đó do máy thứ nhất sản xuất (làm tròn kết quả đến hàng phần trăm).
	\shortans{0,61} 
	\loigiai{ 
		Gọi $A$ là biến cố \lq\lq Chi tiết lấy từ  dây chuyền đạt tiêu chuẩn\rq\rq.\\
		$B$ là  biến cố \lq\lq Chi tiết do máy thứ nhất sản xuất\rq\rq.\\
		Khi  đó $\mathrm{P}(B)=0{,}6$; $\mathrm{P}\big(A \mid B\big)=0{,}9$; $\mathrm{P}\big(A \mid \overline{B}\big)=0{,}85$.\\
		Theo công thức Bayes,  ta có
		 \[
		\mathrm{P}\big(B \mid  A\big)
		= \dfrac{\mathrm{P}\big(B\big)\mathrm{P}\big(A \mid B\big)} 
		{\mathrm{P}\big(B\big)\mathrm{P}\big(A  \mid B\big) + \mathrm{P}\big(\overline{B}\big)\mathrm{P}\big(A \mid \overline{B}\big)}
		= \dfrac{0{,}6 \cdot 0{,}9}{0{,}6  \cdot 0{,}9 + 0{,}4 \cdot 0{,}85}
		=  \dfrac{27}{44} \approx 0{,}61.
		 \]
	 }
\end{bt}

% nguồn: Lop_12_Dai_so_Chuong_6
\begin{bt}%[2D6H2-3]%[ma:4bf643eb]
	An và Bình thi đấu với nhau một trận bóng bàn, người nào thắng trước 3 séc sẽ giành chiến thắng chung cuộc. Xác suất An thắng mỗi séc là $0{,}4$ (không có hòa). Tính xác suất An thắng  chung cuộc. (kết quả làm tròn đến hàng phần trăm)\shortans{0,32}
	 \loigiai{
		Biết số trận thắng trước là $3$ sẽ giành chiến thắng chung cuộc nên An phải chơi tối đa $5$  trận.\\
		Gọi $H$ là  biến cố: \lq\lq An thắng chung cuộc \rq\rq.\\
		 $A$ là biến cố: \lq\lq An thắng trận sau 3 séc \rq\rq.\\
		$B$  là biến cố: \lq\lq An thắng trận sau 4 séc \rq\rq.\\
		$C$ là biến cố: \lq\lq An thắng trận sau 5 séc  \rq\rq .\\
		Khi  đó, $H=A\cup B\cup C$ và $A$, $B$, $C$ đôi một xung khắc.\\
		Áp dụng quy tắc cộng xác  suất, ta có
		$$\mathrm{P}(H)=\mathrm{P}(A)+\mathrm{P}(B)+\mathrm{P}(C).$$ 
		Xác suất An thắng trận sau 3 séc  là $\mathrm{P}(A)=(0{,}4)^3=0{,}064$.\\
		Để An thắng trận sau 4 séc thì An phải  thua 1 trong 3 séc đầu và séc cuối thắng.\\
		Suy ra xác suất An thắng trận sau  4 séc là
		$$\mathrm{P}(C)=\mathrm{C}_4^2\cdot(0{,}4)^3\cdot(0{,}6)^2=0{,}13824$$ 
		Vậy xác suất An thắng chung cuộc  là $\mathrm{P}(H)=0{,}31744\approx0{,}32$.
	} 
\end{bt}

% nguồn: Lop_12_Dai_so_Chuong_6
\begin{bt}%[2D6H2-3]%[ma:99ead346]
	Cho hai biến cố $M$ và $N$, biết rằng $\mathrm{P}(N) = 0{,}7$; $\mathrm{P}(M \mid N) = 0{,}8$ và $\mathrm{P}(M  \mid \overline{N}) = 0{,}4$. Tính $\mathrm{P}(N \mid M)$ (kết quả làm tròn đến hàng phần trăm).
	 \shortans{0,82}
	\loigiai{ 
		Ta có $\mathrm{P}\left(\overline{N}\right)=1-\mathrm{P}(N)=1-0{,}7=0{,}3$.\\ 
		Theo công thức xác suất toàn phần, ta  có
		\[\mathrm{P}(M)=\mathrm{P}(N)\cdot  \mathrm{P}(M\mid N)+\mathrm{P}\left(\overline{N}\right)\cdot\mathrm{P}\left(M\mid\overline{N}\right)=0{,}7\cdot 0{,}8+0{,}3\cdot 0{,}4=0{,}68 \]
		Theo công  thức Bayes, ta có
		\[\mathrm{P}(N\mid  M)=\dfrac{\mathrm{P}(N)\cdot\mathrm{P}(M\mid N)}{\mathrm{P}(M)}=\dfrac{0{,}7\cdot 0{,}8}{0{,}68}=\dfrac{14}{17} \approx 0{,}82. \]
	} 
\end{bt}

% nguồn: Lop_12_Dai_so_Chuong_6
\begin{bt}%[2D6H2-3]%[ma:45a81c40]
	Cho hai biến cố ngẫu nhiên $A$ và $B$. Biết rằng $\mathrm{P}\left(A\mid B\right)=2\mathrm{P}\left(B\mid A\right)$ và  $\mathrm{P}\left(AB\right) \neq 0$. Tính tỉ số $\dfrac{\mathrm{P}\left(A\right)}{\mathrm{P}(B)}$.
	\shortans{2} 
	\loigiai{Theo công thức Bayes, ta có \[\mathrm{P}\left(A\mid B\right)=\dfrac{\mathrm{P}\left(A\right)\cdot\mathrm{P}\left(B\mid A\right)}{\mathrm{P}\left(B\right)} \Rightarrow \dfrac{\mathrm{P}\left(A\right)}{\mathrm{P}\left(B\right)}=\dfrac{\mathrm{P}\left(A\mid  B\right)}{\mathrm{P}(B\mid A)}=2.\]}
\end{bt}

% nguồn: Lop_12_Dai_so_Chuong_6
\begin{bt}%[2D6H2-3]%[ma:4e72cd93]
	Cho hai biến cố $A$, $B$ sao cho $\mathrm{P}(A)=0{,}8$; $\mathrm{P}(B)=0{,}4$ và  $\mathrm{P}(A\mid B)=0{,}9$. Tính $\mathrm{P}(B\mid A)$.
	 \shortans{0,45}
	\loigiai{ 
		Áp dụng công thức Bayes, ta có 
		\[ 
		\mathrm{P}(B\mid A) =  \dfrac{\mathrm{P}(B) \cdot \mathrm{P}(A\mid B)}{\mathrm{P}(A)} = \dfrac{0{,}4 \cdot 0{,}9}{0{,}8} = 0{,}45.
		 \]
	 }
\end{bt}

% nguồn: Lop_12_Dai_so_Chuong_6
\begin{bt}%[2D6H2-3]%[ma:85de5787]
	Một chiếc hộp có $ 50$ viên bi, trong đó có $30$ viên bi màu xanh và $20$ viên bi màu đỏ, các viên bi có kích  thước và khối lượng giống nhau.
	Sau khi kiểm tra, người ta thấy có $70\%$ số viên bi màu xanh được đánh số và $60\%$ số viên bi màu đỏ được đánh số, những viên bi còn lại không đánh số. Lấy ngẫu nhiên một viên bi từ hộp đó. Biết rằng, viên bi lấy ra được đánh số, xác suất để viên bi đó có màu xanh bằng bao  nhiêu (kết quả làm tròn đến hàng phần trăm)? \shortans{0,64}
	 \loigiai{
		Gọi $A$ là  biến cố \lq\lq Viên bi lấy ra có màu xanh\rq\rq.\\
		 Gọi
		$B$ là biến cố \lq\lq Viên bi lấy ra được đánh số\rq\rq.\\ 
		Theo bài ra ta có $\mathrm{P}(A)=\dfrac{30}{50}=\dfrac{3}{5}$; $\mathrm{P}\left(\overline{A}\right)=\dfrac{2}{5}$; $\mathrm{P}(B\mid A)=0{,}7$, $\mathrm{P}\left(B\mid \overline{A}\right)=0{,}6$.\\ 
		Vậy $\mathrm{P}(A\mid B)=\dfrac{\mathrm{P}(A)\cdot \mathrm{P}(B\mid A)}{\mathrm{P}(A)\cdot  \mathrm{P}(B\mid A)+\mathrm{P}\left(\overline{A}\right)\cdot\mathrm{P}\left(B\mid \overline{A}\right)}=\dfrac{\dfrac{3}{5}\cdot 0{,}7}{\dfrac{3}{5}\cdot 0{,}7+\dfrac{2}{5}\cdot 0{,}6}=\dfrac{7}{11}\approx 0{,}64 $.
	 }
\end{bt}

% nguồn: Lop_12_Dai_so_Chuong_6
\begin{bt}%[2D6H2-3]%[ma:811dd69b]
	Tổng kết năm học $2024-2025$, đội HSG toán của CLB chuyên Gia Lai có $7$ bạn được khen thưởng: Phát, Phong, Đức, Kiên, Dương, Khoa và Hải. Phần thưởng cho tất cả các bạn gồm có $4$ quyển sách Đa Thức, $5$ quyển sách Tổ Hợp, $5$ quyển sách Hình Học (các quyển sách cùng chủ đề là giống nhau) sao cho mỗi học sinh được $2$ quyển sách khác chủ đề. Tính xác suất để bạn Khoa và Dương có phần thưởng giống nhau (làm tròn kết quá đến hàng phần  trăm). \shortans{0,24}
	\loigiai{ 
		Gọi $x$ là số học sinh nhận  sách Đa Thức và Tổ Hợp.\\
		Gọi  $y$ là số học sinh nhận sách Đa Thức và Hình Học.\\
		Gọi $z$ là số học sinh nhận  sách Hình Học và Tổ Hợp.\\
		Ta có hệ $\heva{&x+y=4\\&x+z=5\\&y+z=5}\Leftrightarrow\heva{&x=2\\&y=2\\&z=3}$.\\ 
		Không gian  mẫu là $n_{\Omega}=\mathrm{C}_7^2\cdot\mathrm{C}_5^2\cdot\mathrm{C}_3^3=210$.\\
		Gọi $A$ là biến cố "Khoa và Dương  có phần thưởng giống nhau".\\
		Số phần tử của $A$ là $n_A=1\cdot\mathrm{C}_5^2\cdot\mathrm{C}_3^3+1\cdot\mathrm{C}_5^2\cdot\mathrm{C}_3^3+\mathrm{C}_1^1\cdot\mathrm{C}_5^2\cdot\mathrm{C}_3^2=50$.\\ 
		Xác suất để Khoa  và Dương có phần thưởng giống nhau là $\mathrm{P}_A=\dfrac{n_A}{n_{\Omega}}=\dfrac{50}{210}\approx 0{,}24$.
	} 
\end{bt}

% nguồn: Lop_12_Dai_so_Chuong_6
\begin{bt}%[2D6H2-3]%[ma:2c8128b1]
	Bạn An đi học mỗi ngày bằng một trong  hai phương tiện: Xe bus hoặc xe đạp. Vì vội, An chọn ngẫu nhiên một trong hai phương tiện này với xác suất như nhau (tức là $50\%$ đi xe buýt, $50\%$ đi xe đạp). Nếu An đi xe buýt thì xác suất bị muộn học là $6\%$; nếu An đi xe đạp thì xác suất bị muộn học là $4\%$. Hỏi vào một ngày bất kỳ, xác suất An bị muộn học là bao nhiêu.\shortans{0,05}
	 \loigiai{
		Gọi biến cố $B$: \lq\lq An đi xe buýt \rq\rq; $\overline{B}$: \lq\lq An đi xe đạp \rq\rq; $M$:  \lq\lq An đi muộn \rq\rq.\\
		 Xác suất An đi xe buýt $\mathrm{P}(B)=0{,}5$.\\
		Xác suất An đi xe đạp  $\mathrm{P}(\overline{B})=0{,}5$.\\
		Do xác suất An đi học muộn khi đi xe buýt là $6\%$  nên $\mathrm{P}(M\mid B)=0{,}06$.\\
		Do xác suất An đi học muộn khi  đi xe đạp là $4\%$ nên $\mathrm{P}(M\mid B)=0{,}04$.\\
		Xác suất An  bị muộn học là
		$$\mathrm{P}(M)=\mathrm{P}(B)\cdot\mathrm{P}(M\mid B)+\mathrm{P}(\overline{B})\cdot\mathrm{P}(M\mid  \overline{B})=0{,}5\cdot0{,}06+0{,}5\cdot0{,}04=0{,}05.$$
	 }
\end{bt}

% nguồn: Lop_12_Dai_so_Chuong_6
\begin{bt}%[2D6V2-3]%[ma:008dda55]
	Một huyện đang triển khai chương trình sàng lọc phát hiện bệnh hiểm nghèo $X$ cho người dân. Không có loại xét nghiệm nào cho kết quả chính xác 100\%. Một loại xét nghiệm, ta gọi là xét nghiệm $A$, cho kết quả như sau: Khi một người  bị bệnh X thì tỉ lệ có phản ứng dương tính trong xét nghiệm $A$ là 95\%, còn người không mắc bệnh X thì tỉ lệ có phản ứng âm tính trong xét nghiệm $A$ là 97\%. Biết rằng tỉ lệ người mắc bệnh X của huyện này là 0{,}7\%. Hỏi khi một người ở huyện $X$ có xét nghiệm dương tính với xét nghiệm $A$ thì xác suất để người đó mắc bệnh hiểm nghèo $X$ là bao nhiêu? \textit{(kết quả làm tròn đến hàng phần trăm)}.
	 \shortans{0,18}
	 \loigiai{
		Gọi $E$ là biến cố: ``Người đó mắc bệnh'' và $F$ là biến cố:  ``Xét nghiệm dương tính''. \\
		Theo đề  ta có: $\mathrm{P}(E)=0{,}007$; $\mathrm{P}(F|E)=0{,}95$ và $\mathrm{P}(\overline{F}|\overline{E})=0{,}97$. \\
		 Do đó $\mathrm{P}(\overline{E})=1-0{,}007=0{,}993$; $\mathrm{P}(\overline{F}|E)=1-0{,}95=0{,05}$ và $\mathrm{P}(F|\overline{E})=1-0{,}97=0{,}03$. \\
		Xác suất cần tính là $\mathrm{P}(E|F)$.  \\
		Áp dụng công thức Bayes ta có $\mathrm{P}(E|F)=\dfrac{\mathrm{P}(F|E) \cdot \mathrm{P}(E)}{\mathrm{P}(F)}$. \\ 
		Áp dụng công  thức xác suất toàn phần ta có
		 $$\mathrm{P}(F)=\mathrm{P}(F|E)\cdot\mathrm{P}(E)+\mathrm{P}(F|\overline{E})\cdot\mathrm{P}(\overline{E})
		=0{,}95 \cdot 0{,}007+0{,}03  \cdot 0{,}993
		 =0{,}03644.$$
		Vậy $P(E|F)=\dfrac{0{,}95\cdot0{,}007}{0{,}03644}  \approx 0{,}18$.
	 }
\end{bt}

% nguồn: Lop_12_Dai_so_Chuong_6
\begin{bt}%[2D6V2-3]%[ma:5a45f567]
	Tại địa phương A, người ta tiến hành một đợt kiểm tra diện rộng các con bò để phát hiện một loại bệnh X, không có xét nghiệm nào cho kết quả chính xác $100$\%. Có một loại xét nghiệm, mà ở đây ta gọi là xét nghiệm Z cho kết quả như sau: Xét nghiệm có độ nhạy $84$\% (Độ nhạy là xác suất chọn được một mẫu dương tính biết rằng mẫu bị nhiễm bệnh); xác suất dương tính giả là $8$\% (Dương tính giả là xét nghiệm dương tính nhưng thực tế không bị nhiễm bệnh). Biết rằng tỉ  lệ bò ở địa phương A bị mắc bệnh X là $25$\%. Chọn ngẫu nhiên một con bò địa phương A để xét nghiệm, tính xác suất để chọn được con bò bị nhiễm bệnh, biết rằng con bò dương tính với xét nghiệm (kết quả làm tròn đến hàng phần trăm).
	 \shortans{0,78}
	\loigiai{ 
		Gọi $M$ là biến cố \lq \lq Chọn  được một con bò bị nhiễm bệnh\rq \rq.\\
		$N$ là  biến cố \lq \lq Chọn được một con bò dương tính với xét nghiệm\rq \rq.\\
		Theo giả thiết bài toán,  ta có $\mathrm{P}(M)=0{,}25$; $\mathrm{P}(N|M)=0{,}84$; $\mathrm{P}(N|\overline{M})=0{,}08$.\\
		Vậy $\mathrm{P}(M|N)=\dfrac{\mathrm{P}(M).\mathrm{P}(N|M)}{\mathrm{P}(M).\mathrm{P}(N|M)+\mathrm{P}(\overline{M}).\mathrm{P}(N|\overline{M})}=\dfrac{0{,}25\cdot 0{,}84}{0{,}25\cdot 0{,}84+0{,}75\cdot  0{,}08}=\dfrac{7}{9}\approx 0{,}78$.
	} 
\end{bt}

% nguồn: Lop_12_Dai_so_Chuong_6
\begin{bt}%[2D6V2-3]%[ma:0949f799]
	Một đội bắn súng gồm $8$ nam và $2$ nữ. Xác suất bắn trúng của xạ thủ nam là $0{,}8$ và của xạ thủ nữ là  $0{,}9$. Chọn ngẫu nhiên một xạ thủ bắn một viên đạn và thấy bắn trúng. Tính xác suất (làm tròn đến hàng phần trăm) để xạ thủ đó là nữ.
	 \shortans{0,22}
	\loigiai{ 
		Gọi $A$ là biến cố \lq\lq xạ thủ được chọn  là nữ\rq\rq;
		$\overline{A}$ là biến cố \lq\lq xạ thủ là nam\rq\rq \, và $B$  là biến cố \lq\lq bắn trúng\rq\rq. \\
		Khi đó, ta  có
		\[\mathrm \mathrm{P}(A)=\dfrac{2}{10}=\dfrac{1}{5};\quad \mathrm{P}(\overline{A})=\dfrac{4}{5}; \quad \mathrm{P}(B\mid A)=0{,}9;\quad \mathrm{P}(B\mid  \overline{A})=0{,}8.\]
		 Theo xác suất toàn phần, ta có
		 \[\mathrm{P}(B)=\dfrac{1}{5}\cdot0{,}9+\dfrac{4}{5}\cdot0{,}8=0{,}82.
		 \]
		Theo công thức Bayes, ta có 
		 \[
		\mathrm{P}(A\mid B)=\dfrac{\mathrm  \mathrm{P}(A)\mathrm{P}(B\mid A)}{\mathrm \mathrm{P}(B)}=\dfrac{\dfrac{1}{5}\cdot0{,}9}{0{,}82}=\dfrac{9}{41}\approx0{,}22.
		 \]
		Vậy xác suất cần tìm là  $0{,}22$.
	 }
\end{bt}

% nguồn: Lop_12_Dai_so_Chuong_6
\begin{bt}%[2D6V2-3]%[ma:eaef61cb]
	Có ba đồng xu được đựng trong một hộp kín. Đồng xu thứ nhất là một đồng xu cân đối với tỷ lệ mặt ngửa và mặt sấp bằng nhau. Đồng xu thứ hai là một đồng xu bị lỗi có khả năng mặt ngửa xuất hiện là $70$\%. Đồng xu thứ ba là một đồng xu hai mặt ngửa (khi tung luôn ra mặt ngửa). Bạn An lấy ngẫu nhiên một đồng xu từ  hộp và tung nó hai lần. Kết quả của hai lần tung cho thấy xuất hiện một lần mặt sấp và một lần mặt ngửa. Tính xác suất để đồng xu bạn đã chọn là đồng xu thứ hai (đồng xu bị lỗi) (Kết quả làm tròn đến hàng phần trăm).
	\shortans{0,46} 
	\loigiai{ 
		Gọi $A_n$ là biến cố  \lq \lq chọn đồng xu thứ $n$\rq \rq, ( với $n=1;2;3$).\\
		$B$ là biến cố \lq \lq tung  hai lần thì thấy xuất hiện một lần mặt sấp và một lần mặt ngửa\rq \rq.\\
		Vì chọn ngẫu  nhiên nên $\mathrm{P}(A_1)=\mathrm{P}(A_2)=\mathrm{P}(A_3)=\dfrac{1}{3}$.\\
		Lấy ngẫu nhiên  một đồng xu tung hai lần được một mặt sấp và một mặt ngửa thì ta có ba trường hợp như sau
		\begin{itemize}
			\item \textbf{Trường hợp 1:} Chọn được  đồng xu thứ nhất là S-N và N-S nên $\mathrm{P}(B|A_1)=2\cdot\left(\dfrac{1}{2}\cdot\dfrac{1}{2}\right)=\dfrac{1}{2}$.
			\item  \textbf{Trường hợp 2:} Chọn được đồng xu thứ hai là S-N và N-S nên ta có
			$$\mathrm{P}(B|A_2)=0{,}7\cdot 0{,}3+0{,}3\cdot 0{,}7=0{,}42.$$ 
			\item \textbf{Trường hợp 3:} Chọn được đồng  xu thứ ba là N-N nên $\mathrm{P}(B|A_3)=0$.
		\end{itemize}
		Áp dụng công thức Bayes ta tính  được xác suất chọn được đồng xu thứ hai là
		$$\mathrm{P}(A_2|B)=\dfrac{\mathrm{P}(B|A_2).\mathrm{P}(A_2)}{\mathrm{P}(B|A_1).\mathrm{P}(A_1)+\mathrm{P}(B|A_2).\mathrm{P}(A_2)+\mathrm{P}(B|A_3).\mathrm{P}(A_3)}=\dfrac{0{,}42\cdot\dfrac{1}{3}}{\dfrac{1}{2}\cdot\dfrac{1}{3}+0{,}42\cdot\dfrac{1}{3}+0\cdot\dfrac{1}{3}}\approx 0{,}46.$$ 
	 }
\end{bt}

% nguồn: Lop_12_Dai_so_Chuong_6
\begin{bt}%[2D6V2-3]%[ma:f582e833]
	Một cơ quan hành chính nhà nước thực hiện  việc tinh giản biên chế thông qua phỏng vấn. Tỷ lệ nhân viên của cơ quan thuộc hai nhóm trình độ: Đại học, Cao đẳng lần lượt là $65$\% và $35$\%. Qua phỏng vấn thì tỷ lệ nhân viên bị tinh giản của nhóm đại học là $10$\%, nhóm cao đẳng là $15$\%. Chọn một nhân viên bất kỳ đã bị tinh giản thì hãy tính xác suất để người này có trình độ đại học (kết quả là một số thập phân nhỏ hơn $1$ đã làm tròn đến hàng phần trăm). \shortans{0,55}
	\loigiai{ 
		Gọi $A$ là  biến cố: \lq \lq nhân viên được chọn thuộc nhóm trình độ đại học\rq \rq.\\
		$B$ là  biến cố: \lq \lq nhân viên được chọn thuộc nhóm trình độ cao đẳng\rq \rq.\\
		$E$ là biến cố:  \lq \lq nhân viên được chọn bị tinh giản\rq \rq.\\
		Ta có $B=\overline{A}$. Ta cần tính $P\left(E\right)$.\\ 
		Theo công thức xác suất toàn  phần, ta có
		 $$P\left(E\right)=\mathrm{P}\left(A\right).\mathrm{P}(E|A)+\mathrm{P}(\overline{A}).\mathrm{P}(E|\overline{A}).$$
		 Theo bài ra ta có $\mathrm{P}\left(A\right)=0{,}65$, $\mathrm{P}(\overline{A})=\mathrm{P}\left(B\right)=0{,}35$.\\
		$\mathrm{P}(E|A)$ là xác suất để nhân viên  bị tinh giản, thuộc nhóm trình độ đại học.\\
		Theo bài ra ta có $\mathrm{P}(E|A)=0{,}1$.\\ 
		$\mathrm{P}(E|\overline{A})$ là xác suất để nhân viên bị tinh giản, thuộc nhóm trình độ cao đẳng.\\ 
		Theo bài ra ta có $\mathrm{P}(E|\overline{A})=0{,}15$.\\ 
		Thay vào ta được $$P\left(E\right)=\mathrm{P}\left(A\right)\cdot  \mathrm{P}(E|A)+\mathrm{P}(\overline{A})\cdot \mathrm{P}(E|\overline{A})=0{,}65\cdot 0{,}1+0{,}35\cdot 0{,}15=0{,}1175.$$
		Ta có xác suất để nhân viên được chọn có trình độ đại học, biết rằng nhân viên này đã bị  tinh giản, chính là xác suất $\mathrm{P}(A|E)$.\\
		Theo công thức Bayes và kết quả ở trên, ta  có
		$$\mathrm{P}(A|E)=\dfrac{\mathrm{P}\left(A\right).\mathrm{P}(E|A)}{P\left(E\right)}=\dfrac{0{,}65\cdot 0{,}1}{0{,}1175}\approx  0{,}55.$$
	} 
\end{bt}

% nguồn: Lop_12_Dai_so_Chuong_6
\begin{bt}%[2D6V2-3]%[ma:a2b004dc]
	Một siêu thị có hai quầy thanh toán. Quầy $1$ phục vụ $60\%$ lượng khách hàng với tỷ lệ khách hàng có thẻ thành viên là $30\%$. Quầy $2$ phục vụ $40\%$ lượng khách hàng với tỷ lệ khách hàng có thẻ thành viên là $50\%$. Chọn ngẫu nhiên một người rời khỏi quầy thanh toán. Biết rằng  khách hàng này có thẻ thành viên, tính xác suất khách hàng đã thanh toán ở quầy $1$. (làm tròn kết quả đến hàng phần trăm).
	\shortans{0,47} 
	 \loigiai{
		Gọi $A$ là biến cố \lq\lq Khách hàng  thanh toán ở quầy $1$\rq\rq. \\
		Suy ra  $\overline{A}$ là biến cố \lq\lq Khách hàng thanh toán ở quầy $2$\rq\rq. \\
		 Và $B$ là biến cố \lq\lq Khách hàng thanh toán có thẻ thành viên\rq\rq. \\
		Theo đề,  ta có $\mathrm{P}\left(A\right)=0{,}6; \mathrm{P}\left(B\mid A\right)=0{,}3; \mathrm{P}\left(\overline{A}\right)=0{,}4; \mathrm{P}\left(B\mid\overline{A}\right)=0{,}5$. \\
		Ta cần  tính $\mathrm{P}\left(A\mid B\right)$. \\
		Áp dụng công thức Bayes, ta được \\ $$\mathrm{P}\left(A\mid B\right)=\dfrac{\mathrm{P}\left(A\right)\cdot\mathrm{P}\left(B\mid A\right)}{\mathrm{P}\left(A\right)\cdot\mathrm{P}\left(B\mid A\right)+\mathrm{P}\left(\overline{A}\right)\cdot\mathrm{P}\left(B\mid  \overline{A}\right)}=\dfrac{0{,}6\cdot 0{,}3}{0{,}6\cdot 0{,}3+0{,}4\cdot 0{,}5}=\dfrac{0,18}{0{,}38}=\dfrac{9}{19}\approx 0{,}47.$$
		Vậy xác suất khách hàng đã thanh toán ở quầy $1$ biết rằng khách hàng  có thẻ thành viên là $0{,}47$.
	 }
\end{bt}

% nguồn: Lop_12_Dai_so_Chuong_6
\begin{bt}%[2D6V2-3]%[ma:f3ad9518]
	Một hộp chứa $12$ viên bi xanh và $6$ viên bi đỏ có cùng kích thước và khối lượng. Bạn Chiến  lấy ngẫu nhiên ra một viên bi từ hộp, xem màu của nó rồi bỏ ra ngoài. Đến lượt bạn Thắng lấy bi với số lượng phụ thuộc vào màu của viên bi mà bạn Chiến đã lấy. Cụ thể như sau
	\begin{itemize}
		\item Nếu viên  bi bạn Chiến lấy ra có màu xanh thì bạn Thắng sẽ lấy ngẫu nhiên ra hai viên bi từ hộp;
		\item Nếu viên bi bạn Chiến lấy ra có màu đỏ thì bạn Thắng sẽ lấy ngẫu nhiên ra  bốn viên bi từ hộp.
	\end{itemize}
	Tính xác suất để bạn Chiến lấy được viên bi màu đỏ, biết rằng trong các viên bi được bạn Thắng  lấy ra có ít nhất một viên bi khác màu với viên bi của bạn Chiến (làm tròn kết quả đến hàng phần trăm).
	\shortans{0,46} 
	 \loigiai{
		Tổng số bi ban đầu là $12$ xanh $\,+ 6$ đỏ $= 18$ (viên bi). \\ 
		 Gọi $A$ là biến cố \lq\lq bạn Chiến lấy được viên bi màu đỏ\rq\rq. \\
		Suy ra $\overline{A}$ là biến cố  \lq\lq bạn Chiến lấy được viên bi màu xanh\rq\rq. \\
		Và $B$ là biến cố \lq\lq bạn Thắng lấy ra có ít nhất một viên bi khác  màu với viên bi của bạn Chiến\rq\rq. \\
		Ta cần tính $\mathrm{P}\left(A\mid B\right)$. \\ 
		 Ta có $\mathrm{P}\left(A\right)=\dfrac{6}{18}=\dfrac{1}{3}; \mathrm{P}\left(\overline{A}\right)=\dfrac{12}{18}=\dfrac{2}{3}$.
		\begin{itemize}
			\item Nếu viên bi bạn  Chiến lấy ra có màu đỏ thì hộp còn lại $17$ viên (gồm $12$ xanh, $5$ đỏ) và khi đó Thắng sẽ lấy $4$ viên bi.\\
			Biến cố đối $\overline{B}\mid A$ là \lq\lq Thắng lấy cả  $4$ viên cùng màu với viên của Chiến (cùng đỏ)\rq\rq. \\ Ta có $\mathrm{P}\left(\overline{B}\mid A\right)=\dfrac{\mathrm{C}_{5}^{4}}{\mathrm{C}_{17}^{4}}=\dfrac{1}{476}$. \\
			 Suy ra $\mathrm{P}\left(B\mid A\right)=1-\dfrac{1}{476}=\dfrac{475}{476}$.
			\item Nếu viên bi bạn Chiến lấy ra có màu xanh  thì hộp còn lại $17$ viên (gồm $11$ xanh, $6$ đỏ) và khi đó Thắng sẽ lấy $2$ viên bi.\\
			Biến cố đối $\overline{B}\mid \overline{A}$ là \lq\lq Thắng lấy cả $2$ viên cùng màu với viên của Chiến (cùng xanh)\rq\rq. \\ Ta có  $\mathrm{P}\left(\overline{B}\mid \overline{A}\right)=\dfrac{\mathrm{C}_{11}^{2}}{\mathrm{C}_{17}^{2}}=\dfrac{55}{136}$. \\
			Suy ra  $\mathrm{P}\left(B\mid \overline{A}\right)=1-\dfrac{55}{136}=\dfrac{81}{136}$.
		\end{itemize}
		Áp dụng công thức Bayes, ta được $$\mathrm{P}\left(A\mid B\right)=\dfrac{\mathrm{P}\left(A\right)\cdot\mathrm{P}\left(B\mid A\right)}{\mathrm{P}\left(A\right)\cdot\mathrm{P}\left(B\mid A\right)+\mathrm{P}\left(\overline{A}\right)\cdot\mathrm{P}\left(B\mid  \overline{A}\right)}=\dfrac{\dfrac{1}{3}\cdot\dfrac{475}{476}}{\dfrac{1}{3}\cdot\dfrac{475}{476}+\dfrac{2}{3}\cdot\dfrac{81}{136}}= \dfrac{475}{1042}\approx 0{,}46.$$
		Vậy xác suất  để bạn Chiến lấy được viên bi màu đỏ, biết rằng trong các viên bi được bạn Thắng lấy ra có ít nhất một viên bi khác màu với viên bi của bạn Chiến là $0{,}46$.
	 }
\end{bt}

% nguồn: Lop_12_Dai_so_Chuong_6
\begin{bt}%[2D6V2-3]%[ma:58b894e4]
	Bạn Thanh có hai hộp đựng bi, các viên bi có cùng kích thước và hình dạng. Hộp thứ nhất có $8$ viên bi trắng và $2$ viên bi đen, hộp  thứ hai có $6$ viên bi trắng và $3$ viên bi đen. Thanh lấy ngẫu nhiên một viên bi từ hộp thứ nhất để vào hộp thứ hai, sau đó lấy ngẫu nhiên một viên bi từ hộp thứ hai. Biết viên bi lấy ra ở hộp thứ hai là viên bi màu trắng. Tính xác suất viên bi đó thuộc hộp thứ nhất (làm tròn đến hàng phần trăm).
	 \shortans{0,12}
	\loigiai{ 
		Gọi  $A$ là biến cố: \lq\lq Lấy được viên bi trắng ở hộp thứ nhất.\rq\rq;\\
		$B$ là biến cố: \lq\lq Lấy được viên bi trắng ở  hộp thứ hai.\rq\rq
		 \\
		$B_1$ là biến cố: \lq\lq Lấy  được viên bi trắng ở hộp thứ hai đến từ hộp thứ nhất.\rq\rq
		 \\
		Theo đề bài, ta có 
		\begin{align*}
			 &\mathrm{P}(A) = \dfrac{8}{10} &&\mathrm{P}(\overline{A}) = \dfrac{2}{10} \\
			&\mathrm{P}(B \mid A) = \dfrac{7}{10} &&\mathrm{P}(B  \mid \overline{A}) = \dfrac{6}{10}
		\end{align*}
		Xác suất viên bi trắng lấy ra từ hộp thứ  hai thuộc hộp thứ nhất là
		$$ 
		\mathrm{P}(B_1 \mid A) =  \dfrac{\dfrac{8}{10} \cdot \dfrac{1}{10}}{\dfrac{8}{10} \cdot \dfrac{7}{10} + \dfrac{2}{10} \cdot \dfrac{6}{10}}
		= 0{,}12. 
		$$ 
	 }
\end{bt}

% nguồn: Lop_12_Dai_so_Chuong_6
\begin{bt}%[2D6V2-3]%[ma:03b444b8]
	Một huyện đang triển khai chương trình sàng lọc phát hiện bệnh hiểm nghèo $X$ cho người dân. Không có loại xét nghiệm nào cho kết quả chính xác 100\%. Một loại xét nghiệm, ta gọi là xét nghiệm $A$, cho kết quả như sau: Khi một người bị bệnh X thì tỉ lệ có phản ứng dương tính trong xét nghiệm $A$ là 95\%, còn người không mắc bệnh X thì tỉ lệ có phản ứng âm tính trong xét nghiệm $A$ là 97\%. Biết rằng tỉ lệ người mắc bệnh X của huyện này là 0{,}7\%. Hỏi khi một người ở huyện  $X$ có xét nghiệm dương tính với xét nghiệm $A$ thì xác suất để người đó mắc bệnh hiểm nghèo $X$ là bao nhiêu \textit{(kết quả làm tròn đến hàng phần trăm)}?
	 \shortans{0,18}
	\loigiai{ 
		Gọi $E$ là biến cố:  \lq\lq Người đó mắc bệnh\rq\rq\, và $F$ là biến cố: \lq\lq Xét nghiệm dương tính\rq\rq. \\
		Theo đề ta có $\mathrm{P}(E)=0{,}007$;  $\mathrm{P}(F\mid E)=0{,}95$ và $\mathrm{P}\left(\overline{F}\mid \overline{E}\right)=0{,}97$. \\
		 Do đó $\mathrm{P}\left(\overline{E}\right)=1-0{,}007=0{,}993$; $\mathrm{P}(\overline{F}\mid E)=1-0{,}95=0{,05}$ và $\mathrm{P}\left(F\mid \overline{E}\right)=1-0{,}97=0{,}03$. \\
		 Xác suất cần tính là $\mathrm{P}(E\mid F)$. \\
		Áp dụng công thức Bayes ta có $\mathrm{P}(E\mid F)=\dfrac{\mathrm{P}(F\mid E)  \cdot \mathrm{P}(E)}{\mathrm{P}(F)}$. \\
		 Áp dụng công thức xác suất toàn phần ta có
		 $$\mathrm{P}(F)=\mathrm{P}(F\mid E)\cdot\mathrm{P}(E)+\mathrm{P}\left(F\mid \overline{E}\right)\cdot\mathrm{P}\left(\overline{E}\right)
		=0{,}95  \cdot 0{,}007+0{,}03 \cdot 0{,}993
		=0{,}03644.$$ 
		Vậy $P(E\mid  F)=\dfrac{0{,}95\cdot0{,}007}{0{,}03644} \approx 0{,}18$.
	} 
\end{bt}

% nguồn: Lop_12_Dai_so_Chuong_6
\begin{bt}%[2D6V2-3]%[ma:11f342b6]
	Một  công ty có ba phân xưởng sản xuất $A$, $B$, $C$. Xưởng $A$ sản xuất $50\%$ tổng số sản phẩm, xác suất sản phẩm bị lỗi là $2\%$. Xưởng $B$ sản xuất $20\%$ tổng số sản phẩm và xác suất sản phẩm bị lỗi là $3\%$. Xưởng $C$ sản xuất $30\%$ tổng số sản phẩm, xác suất sản phẩm bị lỗi là $5\%$. Chọn ngẫu nhiên một sản phẩm bị lỗi, hỏi xác suất sản phẩm đó do xưởng $B$ sản xuất là bao nhiêu? (Kết quả làm tròn đến phần trăm).
	 \shortans{0,19}
	 \loigiai{
		 Gọi
		\begin{itemize}
			\item $A$, $B$, $C\colon $ \lq\lq biến cố sản phẩm được sản xuất từ xưởng  $A$, $B$, $C$\rq\rq.
			\item  $L\colon $ \lq\lq biến cố sản phẩm bị lỗi\rq\rq.
		\end{itemize}
		Theo đề bài 
		\begin{align*}
			\mathrm{P}(A) &= 0{,}5,\quad \mathrm{P}(B) =  0{,}2,\quad \mathrm{P}(C) = 0{,}3,\\
			\mathrm{P}(L|A) &=  0{,}02,\quad \mathrm{P}(L|B) = 0{,}03,\quad \mathrm{P}(L|C) = 0{,}05.
		\end{align*}
		Áp dụng định lý xác suất toàn phần 
		\begin{eqnarray*}
			&\mathrm{P}(L)&= \mathrm{P}(A)\mathrm{P}(L|A) + \mathrm{P}(B)\mathrm{P}(L|B) + \mathrm{P}(C)\mathrm{P}(L|C)\\ 
			&&= 0{,}5  \cdot 0{,}02 + 0{,}2 \cdot 0{,}03 + 0{,}3 \cdot 0{,}05\\
			&&= 0{,}01 + 0{,}006 +  0{,}015 = 0{,}031.
		\end{eqnarray*}
		Áp dụng công thức Bayes để tính $\mathrm{P}(B|L)$ 
		 \[
		\mathrm{P}(B|L)  = \dfrac{\mathrm{P}(B)\mathrm{P}(L|B)}{\mathrm{P}(L)} = \dfrac{0{,}2 \cdot 0{,}03}{0{,}031} = \dfrac{0{,}006}{0{,}031} \approx 0{,}19.
		 \]
	} 
\end{bt}

% nguồn: Lop_12_Dai_so_Chuong_6
\begin{bt}%[2D6V2-3]%[ma:4a1a645e]
	 Trường THPT Thị xã Quảng trị, có $20\%$ học sinh tham gia câu lạc bộ âm nhạc, trong số học sinh đó thì có $75\%$ học sinh biết chơi đàn guitar. Ngoài ra, trong số học sinh không tham gia câu lạc bộ âm nhạc có $10\%$ học sinh biết chơi đàn guitar. Chọn ngẫu nhiên $1$ học sinh của trường. Giả sử học sinh đó biết chơi đàn guitar. Xác suất để chọn được học sinh thuộc câu lạc bộ âm nhạc là bao nhiêu? (làm tròn đến hàng phần trăm).
	 \shortans{0,65}
	 \loigiai{
		Gọi $A$ là biến cố học sinh tham gia CLB âm nhạc, ta  có $\mathrm{P}(A)=0{,}2\Rightarrow \mathrm{P}(\overline{A})=1-0{,}2=0{,}8$.\\
		Gọi $G$ là biến cố học sinh biết chơi đàn guitar, ta có $\mathrm{P}(G\mid A) = 0{,}75$;  $\mathrm{P}(G\mid \overline{A}) = 0{,}1$.\\
		Xác suất để chọn được học sinh thuộc CLB âm nhạc, biết rằng học sinh đó biết chơi đàn  guitar là $$\mathrm{P}(A\mid G)=\dfrac{\mathrm{P}(G\mid A)\cdot \mathrm{P}(A)}{\mathrm{P}(G\mid A)\cdot \mathrm{P}(A)+\mathrm{P}(G\mid\overline{A})\cdot \mathrm{P}(\overline{A})}=\dfrac{0{,}75 \cdot 0{,}2}{0{,}75 \cdot 0{,}2 + 0{,}1 \cdot 0{,}8 } = \dfrac{0{,}15}{0{,}15 + 0{,}08} \approx 0{,}65.$$
	} 
\end{bt}

% nguồn: Lop_12_Dai_so_Chuong_6
\begin{bt}%[2D6V2-3]%[ma:62af95cb]
	Một bệnh viện sử  dụng một xét nghiệm chẩn đoán bệnh $X$. Xét nghiệm này có độ chính xác $95 \%$ khi bệnh nhân thực sự mắc bệnh $X$ (tức là, nếu bệnh nhân mắc bệnh $X$, xét nghiệm sẽ cho kết quả dương tính với xác suất $95\%$). Tuy nhiên, xét nghiệm trện cũng có thể cho kết quả dương tính giả (false positive) với xác suất $5\%$ ở những người không mắc bệnh $X$. Giả sử tỉ lệ mắc bệnh $X$ trong dân số là $1\%$. Khi một bệnh nhân được xét nghiệm và cho kết quả dương tính. Tính xác suất bệnh nhân này thực sự mắc bệnh $X$ (kết quả được ghi dưới dạng số thập phân, làm tròn đến hàng phần trăm).
	 \shortans{0,16}
	 \loigiai{
		Gọi $A$ là biến cố \lq\lq bệnh nhân mắc bệnh $X$\rq\rq, $B$ là biến cố \lq\lq xét nghiệm cho kết quả dương  tính\rq\rq.\\
		Ta có $\mathrm{P}(AB)=\mathrm{P}(B|A)\cdot \mathrm{P}(A)=0{,}95\cdot 0{,}01=0{,}0095$.\\ 
		$\mathrm{P}(B)=\mathrm{P}(B| A)\cdot \mathrm{P}(A)+\mathrm{P}(B|  \overline{A})\cdot \mathrm{P}(\overline{A})=0{,}95\cdot 0{,}01+0{,}05\cdot 0{,}99=0{,}059$.\\
		Xác suất  bệnh nhân này thực sự mắc bệnh X là $\mathrm{P}(A| B)=\dfrac{\mathrm{P}(AB)}{\mathrm{P}(B)}=\dfrac{0{,}0095}{0{,}055}=\dfrac{19}{118}\approx 0{,}16$.
	} 
\end{bt}

% nguồn: Lop_12_Dai_so_Chuong_6
\begin{bt}%[2D6V2-3]%[ma:979aaa6d]
	Ở một viện dưỡng lão, tỉ lệ người mắc bệnh tim mạch là $25\%$. Tỉ lệ người hút thuốc lá trong số những người mắc bệnh tim mạch là một số dương và bằng $2$ lần tỉ lệ người hút thuốc lá trong số những người không mắc bệnh tim mạch. Hỏi xác  suất để một người ở viện dưỡng lão này mắc bệnh tim mạch là bao nhiêu phần trăm, biết rằng người đó hút thuốc?
	\shortans{40} 
	 \loigiai{
		Gọi  các biến cố như sau
		\begin{itemize}
			\item $A$ là  biến cố \lq\lq Người ở viện dưỡng lão mắc bệnh tim mạch\rq\rq;
			\item $B$ là biến cố \lq\lq Người ở viện dưỡng lão  hút thuốc\rq\rq.
		\end{itemize}
		Khi  đó xác suất cần tính là $\mathrm{P}\left(A\mid B\right)$. \\
		Ta có $\mathrm{P}\left(A\right)=25\%=0{,}25\Rightarrow  \mathrm{P}\left(\overline{A}\right)=1-\mathrm{P}\left(A\right)=0{,}75$. \\
		Và $\mathrm{P}\left(B\mid  A\right)=2\mathrm{P}\left(B\mid \overline{A}\right)=2p, \text{ với } \mathrm{P}\left(B\mid \overline{A}\right)=p, (p>0)$. \\
		Áp dụng công  thức Bayes, ta được
		$$\mathrm{P}\left(A\mid B\right)=\dfrac{\mathrm{P}\left(A\right)\cdot\mathrm{P}\left(B\mid A\right)}{\mathrm{P}\left(A\right)\cdot\mathrm{P}\left(B\mid A\right)+\mathrm{P}\left(\overline{A}\right)\cdot\mathrm{P}\left(B\mid \overline{A}\right)}=\dfrac{0{,}25\cdot 2p}{0{,}25\cdot 2p+0{,}75\cdot p}=\dfrac{0{,}25\cdot  2}{0{,}25\cdot 2+0{,}75}=\dfrac{0{,}5}{1{,}25}=\dfrac{2}{5}=40\%.$$
		Vậy  xác suất để một người ở viện dưỡng lão này mắc bệnh tim mạch, biết rằng người đó hút thuốc là $40\%$.
	 }
\end{bt}

% nguồn: Lop_12_Dai_so_Chuong_6
\begin{bt}%[2D6V2-3]%[ma:a5863646]
	Một phân xưởng sản xuất bóng đèn có tỉ lệ bóng đạt chuẩn là $95 \%$. Để hạn chế số lượng bóng  không đạt chuẩn được bán ra thị trường, người ta lắp đặt một thiết bị kiểm tra chất lượng tự động $\mathrm{S}$. Nếu một bóng đèn không đạt chuẩn, thiết bị $\mathrm{S}$ sẽ loại bỏ nó với xác suất $0{,}99$. Khi kiểm tra lại các bóng đèn bị loại, người ta nhận thấy có $10 \%$ số đó là bóng đạt chuẩn. Chọn ngẫu nhiên $1$ bóng đèn do phân xưởng đó sản xuất. Xác suất bóng đèn được chọn đạt chuẩn biết rằng nó không bị thiết bị $\mathrm{S}$ loại bỏ bằng $\dfrac{a}{b}$ với $a$, $b$ là số nguyên dương và $b<2000$. Giá trị của biểu thức $a+b$ là bao nhiêu?
	\shortans{3779} 
	\loigiai{ 
		Chọn ngẫu nhiên $1$ bóng  đèn do phân xưởng sản xuất.\\
		Gọi $C$ là biến cố bóng đèn  đó đạt chuẩn và $L$ là biến cố bóng đèn đó bị thiết bị $\mathrm{S}$ loại.
		Ta có $\mathrm{P}(C)=0{,}95$; $\mathrm{P}\left(L \mid \overline{C}\right)=0{,}99$;  $\mathrm{P}(C \mid L)=0{,}1$.\\
		Suy ra $\mathrm{P}(L)=\dfrac{\mathrm{P}\left(L \mid \overline{C}\right) \mathrm{P}\left(\overline{C}\right)}{\mathrm{P}\left(\overline{C} \mid L\right)}=\dfrac{0{,}99 \cdot 0{,}05}{1-0{,}1}=\dfrac{11}{200}$\\ 
		và $\mathrm{P}(CL)=\mathrm{P}(C  \mid L) \cdot P(L)=0{,}1 \cdot \dfrac{11}{200}=\dfrac{11}{2000}$.\\
		Xác suất bóng đèn được chọn đạt  chuẩn biết rằng nó không bị thiết bị $\mathrm{S}$ loại là
		 $$\mathrm{P}\left(C \mid \overline{L}\right)=\dfrac{\mathrm{P}\left(\mathrm{C}\overline{L}\right)}{\mathrm{P}\left(\overline{L}\right)}=\dfrac{\mathrm{P}(C)-\mathrm{P}(CL)}{1-\mathrm{P}(L)}=\dfrac{1889}{1890}.$$
		Suy ra  $a+b=3779$.
	 }
\end{bt}

% nguồn: Lop_12_Dai_so_Chuong_6
\begin{bt}%[2D6V2-3]%[ma:18fad900]
	Trong một kì thi tuyển sinh, một trường đại học thống kê rằng $70\%$ thí sinh lựa chọn khối D01 (gồm các môn Toán, Văn, Tiếng Anh). Biết rằng, nếu thí sinh chọn khối D01 thì xác suất để thí sinh đậu đại học là $0{,}75$;  còn nếu thí sinh không chọn khối D01 thì xác suất để thí sinh đậu đại học là $0{,}6$. Chọn ngẫu nhiên một thí sinh đã tham dự vào kì tuyển sinh của trường biết thí sinh đó đã đậu đại học. Tính xác suất để thí sinh đó chọn khối D01 (Kết quả làm tròn đến hàng phần trăm).
	 \shortans{0,74}
	\loigiai{ 
	 Gọi $A$: \lq\lq Thí sinh đó chọn khối D01\rq\rq; $B$: \lq\lq Thí sinh đó đậu đại học\rq\rq.\\
	Áp dụng công thức  Bayes ta được
	$$P(A|B)=\dfrac{P(A)\cdot P(B|A)}{P(A)\cdot P(B|A)+P\left(\overline{A}\right)\cdot P(B|\overline{A})}=\dfrac{0{,}7\cdot 0{,}75}{0{,}7\cdot 0{,}75+0{,}3\cdot 0{},6}\approx  0{,}74.$$
	} 
\end{bt}

% nguồn: Lop_12_Dai_so_Chuong_6
\begin{bt}%[2D6V2-3]%[ma:ef706e8d]
	Email Filter là một phần mềm chặn email quảng cáo. Nếu một email là thư quảng cáo, phần mềm sẽ chuyển nó vào thư mục Spam với xác suất là $0{,}9$. Ngược lại, nếu một email không là thư quảng cáo, phần mềm có thể chuyển nó vào thư mục Spam với xác suất là $0{,}05$. Thống kê trong một số lượng lớn email bị chuyển vào thư mục Spam thì thấy tỉ lệ thư quảng cáo là $72\%$. Xác suất một email là thư quảng  cáo là bao nhiêu? (Làm tròn kết quả đến hàng phần trăm).
	\shortans{0,13} 
	\loigiai{Chọn ngẫu nhiên  một email. Gọi $A$ là biến cố email đó là thư quảng cáo và $B$ là biến cố Email Filter chuyển email đó vào thư mục Spam.\\
		 Ta có $\mathrm{P}(B|A)=0{,}9$; $\mathrm{P}(B|\overline{A})=0{,}05$;
		 $\mathrm{P}(A|B)=0{,}72$.\\
		Áp  dụng công thức Bayes, ta có
		\begin{center}
			$\mathrm{P}(A|B)=\dfrac{\mathrm{P}(B|A)\cdot  \mathrm{P}(A)}{\mathrm{P}(B|A)\cdot \mathrm{P}(A)+\mathrm{P}(B|\overline{A})\cdot \mathrm{P}({\overline{A}})}$.
		\end{center}
		Đặt $\mathrm{P}(A)=p\in  \left[0;1\right]$, ta có
		\begin{center}
			$0{,}72=\dfrac{0{,}9p}{0{,}9p+0{,}05\left(1-p\right)}\Leftrightarrow 0{,}8\left(0{,}85p+0{,}05\right)=p\Leftrightarrow  p=0{,}125\approx 0{,}13$.
	\end{center}}
\end{bt}

% nguồn: Lop_12_Dai_so_Chuong_6
\begin{bt}%[2D6V2-3]%[ma:e434bd91]
	 Trong kì thi tốt nghiệp THPT, trường A có $60\%$ học sinh chọn khối D. Biết rằng nếu chọn khối D thì xác suất đỗ đại học là $0{,}7$, còn nếu không chọn khối D thì xác suất đỗ là $0{,}8$. Chọn ngẫu nhiên một học sinh đã đỗ đại học, xác suất để học sinh đó chọn khối D là $\dfrac{m}{n}$ (tối giản). Tính $m+n$.
	\shortans{58} 
	\loigiai{ 
		Gọi $A$ là biến cố \lq\lq học  sinh chọn khối D\rq\rq; $B$ là biến cố \lq\lq đỗ đại học\rq\rq.\\
		Khi đó,  ta có
		\[\mathrm \mathrm{P}(A)=0{,}6;\quad \mathrm{P}\left(\overline{A}\right)=0{,}4;\quad \mathrm{P}(B\mid A)=0{,}7;\quad \mathrm{P}(B\mid \overline{A})=0{,}8.\] 
		Vậy 
		\[ 
		\mathrm{P}(A\mid B)=\dfrac{0{,}6\cdot0{,}7}{0{,}6\cdot0{,}7+0{,}4\cdot0{,}8}=\dfrac{21}{37}\Rightarrow  m+n=21+37=58.
		 \]
	} 
\end{bt}

% nguồn: Lop_12_Dai_so_Chuong_6
\begin{bt}%[2D6C2-3]%[ma:cb12995b]
	Một nhóm có $5$ học sinh nam và $4$ học sinh nữ tham gia lao động trên sân trường. Cô giáo chọn ngẫu nhiên đồng thời hai bạn trong nhóm đi tưới cây. Tính xác suất để hai bạn được chọn có cùng giới tính, biết rằng  có ít nhất một bạn nam được chọn. (Kết quả làm tròn đến hai chữ số thập phân).
	\shortans{0,33} 
	\loigiai{ 
		Gọi $A$ là  biến cố \lq\lq Hai bạn được chọn có cùng giới tính\rq\rq.\\
		Gọi $ B $ là biến cố \lq\lq Có ít nhất một  bạn nam được chọn\rq\rq.\\
		$ \Rightarrow AB$ là biến  cố \lq\lq Hai bạn được chọn là nam\rq\rq.\\
		Xác suất để chọn được hai bạn nam là $\mathrm{P}\left(AB\right)=\dfrac{\mathrm{C}_5^2}{\mathrm{C}_9^2}=\dfrac{5}{18}$.\\ 
		Xác suất để  chọn được ít nhất 1 bạn nam $ \mathrm{P}\left(B\right)=\dfrac{\mathrm{C}_5^1\cdot\mathrm{C}_4^1+\mathrm{C}_5^2}{\mathrm{C}_9^2}=\dfrac{5}{6}$.\\
	$ \mathrm{P}\left(A\mid  B\right)=\dfrac{\mathrm{P}\left(AB\right)}{\mathrm{P}\left(B\right)}=\dfrac{1}{3}\approx 0{,}33$.}
\end{bt}

% nguồn: Lop_12_Dai_so_Chuong_6
\begin{bt}%[2D6C2-3]%[ma:a1286ad8]
	Tỷ lệ phế phẩm của một công ty là $10\%$. Trước khi đưa ra thị trường, các sản phẩm được kiểm tra bằng máy nhằm loại bỏ phế phẩm. Xác suất để máy nhận biết đúng chính phẩm là $95\%$, nhận biết đúng phế phẩm là $90\%$. Khi đó tỉ lệ phế phẩm của công  ty trên thị trường bằng $\dfrac{a}{b}$, $\dfrac{a}{b}$ là phân số tối giản. Tính $a+b$.\shortans{175}
	\loigiai{\textbf{Tỉ lệ phế  phẩm của công ty trên thị trường bằng tỉ lệ của số phế phẩm được bán ra trên tổng số sản phẩm được bán ra.}\\
		Gọi $e$  là tỉ lệ phế phẩm kết luận đúng,
		$f$ là tỉ  lệ phế phẩm kết luận sai,
		$c$ là tỉ lệ chính  phẩm kết luận đúng,
		 $d$ là tỉ lệ chính phẩm kết luận sai.\\
		Ta có hệ  phương trình sau
		$$\heva{&e+f+c+d=1\\ &e+f=0{,}1 \\ &\dfrac{e}{e+f}=0{,}9 \\  &\dfrac{c}{c+d}=0{,}95 } \Leftrightarrow \heva{&e+f+c+d=1 \\ &e+f=0{,}1 \\ &0{,}1e-0{,}9f=0 \\ &0{,}05c-0{,}95d=0} \Leftrightarrow\heva{&e=0{,}09 \\ &f=0{,}01 \\ &c=0{,}855 \\ &d=0{,}045.}$$
		Vậy tỉ lệ phế phẩm của công ty trên thị trường  là $P_f=\dfrac{f}{f+c}= \dfrac{0{,}01}{0{,}01+0{,}855}=\dfrac{2}{173}$.\\
		Vậy $a+b=2+173=175$. 
	 }
\end{bt}

% nguồn: Lop_12_Dai_so_Chuong_6
\begin{bt}%[2D6C2-3]%[ma:1c62d288]
	Trong cộng đồng, tỉ lệ tự nhiên của các nhóm máu $O$, $A$, $B$, $AB$ lần lượt là $33{,}7\%$, $37{,}5\%$, $20{,}9\%$ và  $7{,}9\%$. Lấy ngẫu nhiên một người cần máu và $1$ người hiến máu. Hỏi xác suất có thể thực hiện truyền máu là bao nhiêu? (Làm tròn đến hàng phần trăm).
	\shortans{0,57} 
	\loigiai{ 
		Gọi $H$ là biến cố có thể thực hiện truyền máu.\\ 
		Gọi $O$ là biến cố người  nhận có nhóm máu $O$. \\
		Khi đó, người hiến chỉ có thể có  nhóm máu $O$.\\
		$ \Rightarrow \mathrm{P}\left(H\mid O\right) = 0{,}337$.\\ 
		Gọi $A$ là biến cố người nhận  có nhóm máu $A$. \\
		Khi đó, người hiến có  thể có nhóm máu $O$ và $A$.\\
		$ \Rightarrow \mathrm{P}\left(H\mid A\right) = 0{,}337 +  0{,}375 $.\\
		 Gọi $B$ là biến cố người nhận có nhóm máu $B$. \\
		Khi đó, người hiến  có thể có nhóm máu $O$ và $B$.\\
		$ \Rightarrow \mathrm{P}\left(H\mid B\right) = 0{,}337  + 0{,}209 $.\\
		Gọi $C$ là biến cố người  nhận có nhóm máu $AB$. \\
		Khi đó, người hiến có thể có nhóm máu $O$, $A$,  $B$ và $AB$.\\
		$ \Rightarrow \mathrm{P}\left(H\mid C\right) = 0{,}337 +  0{,}375 + 0{,}209 + 0{,}079=1$.\allowdisplaybreaks
		\begin{eqnarray*}
			 \mathrm{P}\left(H\right)&=&\mathrm{P}\left(O\right)\cdot\mathrm{P}\left(H\mid O\right)+\mathrm{P}\left(A\right)\cdot\mathrm{P}\left(H\mid A\right)+\mathrm{P}\left(B\right)\cdot\mathrm{P}\left(H\mid B\right)+\mathrm{P}\left(C\right)\cdot\mathrm{P}\left(H\mid C\right) \\
			&=& 0{,}337\cdot0{,}337 + 0{,}375\left(0{,}337+0{,}375\right)+0{,}209\left(0{,}337+0{,}209\right)+0{,}079\cdot1\\ 
			&=& 0{,}573683. 
		\end{eqnarray*}
	Vậy xác suất có thể truyền máu là $0{,}57$.} 
\end{bt}

% nguồn: Lop_12_Dai_so_Chuong_6
\begin{bt}%[2D6C2-3]%[ma:8c4357e6]
	Ba cầu thủ sút phạt đền $11$ m, mỗi người đá một lần với xác suất làm bàn tương ứng  $a$; $b$ và $0{,}7$ (với $0<b<a<1$). Biết xác suất ghi bàn để ít nhất trong ba cầu thủ ghi bàn là $0{,}982$ và xác suất để ba cầu thủ ghi bàn là $0{,}392$. Tính xác suất để có đúng hai cầu thủ ghi bàn (làm tròn kết quả đến hàng phần trăm).
	\shortans{0,43} 
	\loigiai{ 
		Gọi $A_i $  là biến cố \lq\lq người thứ $ i $ ghi bàn\rq\rq\, với $i=\overline{1,3}$.\\
		Ta có các biến cố $A_1$, $A_2$, $A_3$ là các biến  cố độc lập và $\mathrm{P}\left(A_1\right)=a$, $\mathrm{P}\left(A_2\right)=b$, $\mathrm{P}\left(A_3\right)=0{,}7$.\\
		Gọi $A$ là biến cố \lq\lq Có ít nhất  một trong ba cầu thủ ghi bàn\rq\rq;
		$B$ là biến cố \lq\lq  Cả ba cầu thủ đều ghi bàn\rq\rq;
		$C$ là  biến cố \lq\lq Có đúng hai cầu thủ ghi bàn\rq\rq.\\
		Ta  có
		$\overline{A}=\overline{A_1}\overline{A_2}\overline{A_3} \Rightarrow \mathrm{P}\left(\overline{A}\right) = \mathrm{P}\left(\overline{A_1}\right)  \cdot \mathrm{P}\left(\overline{A_2}\right) \cdot \mathrm{P}\left(\overline{A_3}\right) = 0{,}3 \cdot \left(1-a\right) \cdot \left(1-b\right) $.\\
		 $ \Rightarrow \mathrm{P}\left(A\right)=1-0{,}3 \cdot\left(1-a\right)\cdot\left(1-b\right) $.\\
		Lại có  $B=A_1A_2A_3 \Rightarrow \mathrm{P}\left(B\right)=\mathrm{P}\left(A_1\right)\cdot\mathrm{P}\left(A_2\right)\cdot\mathrm{P}\left(A_3\right) = 0{,}7ab$.\\
		Từ giả  thiết ta có hệ phương trình
		$$\heva{&1-0{,}3\cdot\left(1-a\right)\cdot\left(1-b\right)=0{,}982 \\ &0{,}7ab=0{,}392} \Leftrightarrow \heva{&1+ab-\left(a+b\right)=0{,}06 \\ &ab=0{,}56} \Leftrightarrow \heva{&a+b=1{,}5 \\&ab=0{,}56} \Rightarrow \heva{&a=0{,}8\\&b=0{,}7} \left( \text{ do }  a >b\right).$$
		Mặt  khác ta có $C=\overline{A_1}A_2A_3+A_1\overline{A_2}A_3+A_1A_2\overline{A_3}$ nên
		$$\mathrm{P}\left(C\right)=\left(1-a\right)\cdot b\cdot0{,}7+a\cdot\left(1-b\right)\cdot0{,}7+a\cdot b\cdot0{,}3=0{,}2\cdot0{,}7\cdot0{,}7+0{,}8\cdot0{,}3\cdot0{,}7+0{,}8\cdot0{,}7\cdot0{,}3 = 0{,}434.$$ 
	} 
\end{bt}

% nguồn: Lop_12_Dai_so_Chuong_6
\begin{bt}%[2D6C2-3]%[ma:4c5031c8]
	Có $40$ phiếu thi Toán $12$, mỗi phiếu chỉ có một câu hỏi, trong đó có $13$ câu hỏi lý thuyết (gồm $5$ câu hỏi khó và $8$ câu hỏi dễ) và $27$ câu hỏi bài tập (gồm $12$  câu hỏi khó và $15$ câu hỏi dễ). Lấy ngẫu nhiên ra một phiếu. Tìm xác suất rút được câu hỏi lý thuyết, biết rằng đó là câu hỏi khó. (Kết quả làm tròn đến hàng phần trăm).
	\shortans{0,29} 
	\loigiai{ 
		Gọi $A$ là biến cố \lq\lq  Rút ra được câu hỏi lý thuyết\rq\rq.\\
		Gọi $B$  là biến cố \lq\lq Rút ra được câu hỏi khó\rq\rq.\\
		Nếu biến $B$ đã xảy ra (nghĩa là câu hỏi rút ra là một trong số $17$ câu khó) thì xác suất để câu hỏi đó là lý thuyết (nghĩa là câu hỏi đó là một trong số $5$ câu hỏi lý thuyết khó) chính là xác suất $A$ có điều  kiện $B$ đã xảy ra. Ta đi tính $\mathrm{P}\left(A\mid B\right)$.\\
		$ \mathrm{P}\left(A\right)=\dfrac{13}{40}$, $\mathrm{P}\left(B\right)=\dfrac{17}{40}$, 
	$ \mathrm{P}\left(A \cap B\right)=\dfrac{5}{40}$, $\mathrm{P}\left(A\mid  B\right)=\dfrac{\tfrac{5}{40}}{\tfrac{17}{40}}=\dfrac{5}{17}\approx 0{,}29$.}
\end{bt}
